%\documentclass[11pt,final,openany,mytheorems]{memoir}
\documentclass[11pt]{book}
% \documentclass{ar-1col}
% \usepackage[numbers]{natbib}
% \usepackage{url}
\input{preamble}
\usepackage{fullpage}
\begin{document}
%\frontmatter
\title{{\bf \large Project Kuiper Synchronization and Positioning}}
% {\small {\color{magenta} {\bf Disclaimer:} Notes for group discussion only ... \\[-.1in]
% will be polished overtime for a wider distribution}}}
\author{\large Mehran Mesbahi}
\date{\color{blue} \bf February 26, 2024}
%\makehalftitle
% \begin{figure}
% \begin{center}
% \scalebox{1.}{\includegraphics{eps/download.jpg}}
% \end{center}
% \end{figure}
% %

% \newpage
% % \thispagestyle{empty}

% \vspace{1in}

% \begin{shaded*}
% \section*{{\color{blue} Acknowledgments}}
% Weekly discussions
% with Tim Gallagher, Pengfei Xia, Vikram Chandrasekhar, Hongbing Cheng, Iskender Kushan, Duong Hoang, and Tim Deck, 
% instrumental in motivating the setup and scenarios
% analyzed in this report, are much appreciated. 
% Furthermore, the Kuiper team's sharing of the Matlab code for Kuiper synchronization
% as well as the ``confluence'' page is acknowledged.
% \end{shaded*}
% \newpage
%------------------------
% \vspace{0.01in}
% \begin{shaded*}
% \section*{{\color{cyan} Current work as of 5/16/2022}}
% \begin{itemize}
% 	\item {\color{blue} Two-way time transfer}
% 	\item {\color{blue} Kuiper aerial vehicle network}
% 	\item  {\color{blue} Using the KAV onboard camera for (initial) position/orientation estimation of Kuiper ground rover}
% 	\item  {\color{blue} IMU/Kuiper integration schemes for mobile units- documentation}
% 	\item  {\color{blue} Aerial vehicles, IMU/Kuiper integration}\\[.1in]
	
%     \hline \hline
% 	\item Use of inter-vehicle message passing for differential positioning and timing-- this is very important area of current interest and lot of applications for mobile units, as well as indoor units- please advise on priority.
% 	\item Intermittent filtering (effect of interruptions on SV signal reception) for simultaneous synchronization and localization; performance characterization; this is more of a research question- how much intermittency can be tolerated in the synchronization and timing protocol?
% 	\item {Geometric dilution of precision for mobile units with Kuiper SV time stamping (Chapter 12)- theory done; advise on priority}
% 	\item {Synchronization with other sets of observables, carrier beat frequency; nav using comm; advise on priority}
% 	\item Complete the analysis on ISL synchronization; network coherence and network bias (Chapter 10; done)
% 	\item how to go from characterizing ``important'' orbits for synchronization to GW site location-- optimizing over ground tracks; done for the most part-- Chapter 10
% 	\item Differencing schemes and applications to Kuiper (Chapter 5)
% 	\item Refinements on the CT receiver algorithms- weighted least squares as a function of the path delay covariances and location of SVs with respect to CTs	
% 	\item IoT synchronization; synchronization/localization for indoor units
% 	\item Kuiper/GPS integration for higher performance synchronization and timing
	
% \end{itemize}
% \end{shaded*}

% \newpage
% \thispagestyle{empty}
% \begin{shaded*}
% {\color{red} \bf Current Work:}

% \begin{itemize}
% \item distributed synchronization; synchronization on standalone Kuiper network/basic architecture; performance measures
% \item Kalman filtering for clock synchronization; ensemble Kalman filtering, clock ensembles
% \item MLE of clock parameters with respect to delay statistics (exponential and Gaussian)
% \item normal GPS operation; with SV and CT
% \item iterative receiver design/factor graphs (chapter 8); synchronization protocols using factor graphs
% \end{itemize}
% \end{shaded*}

\newpage
\tableofcontents
% \input{preface}
% \input{preface2}
% \newpage
% \input{notation}
%\mainmatter
%------------------------------


% {\color{red} \bf art work on Kuiper synchronization}


% \begin{itemize}

% \item learning the clock model; bias, skew, and covariance? or do it in one shot?
% \item delay propagation estimate
% \item learning the delays 
% \item two state model of clock and Allan variance- Kalman filter parameters
% \item the physics of synchronization
% \item inter-satellite clock synchronization
% \item review the code by Kuiper team
% }
% \end{itemize}
% \include{notation}

% \chapter{foundational aspects}

% what are the fundamental issues in constellations synchronization?
\part{2023 Effort}
\include{gps-kuiper}
\include{sync-ISL-I}
% \include{sync-ISL-II}


# Code Log: November 27, 2023; Mehran Mesbahi
# Dec. 26, 2023
# new changes in the Kuper CT class and a few new functionalities
# added the option to correct for the clock recursively
# added velocity state for the KCTs
# added doppler measurement scheme
# added doppler based extended Kalman filter capability 
# combined measurement schemes for EKF based PNT
# -----------
# Nov. 8, 2023
# - most of the changes has been implemented in class Kuiper_CT 
# - added a fisher information metric to choose which satellite to track to reduce position uncertainty 
# - added an incentive cost not to do the handover for a marginal incremental PNT benefit 
# - a number of fixes on scaling, particularly for ``speedOflight'' parameter (m/s vs. km/s) to be consistent 
# - fixed some indexing issues with sate llites identified to be tracked by KCT
# - fixed a typo in the extended Kalman filter update
# - included an iterative options for the extended Kalman filtering
# - included a graphical means of visualizing the prediction from the extended KF and the actual estimation error


\end{document}
%----
\include{2023-topics}
\part{}
\include{coordination-synchronization}
\include{summary-report}

%\part{Preliminaries}
\include{introduction}
%\part{Building Blocks}
%\include{intro-to-synchronization}
\include{preliminaries}

% \include{coordinate-frames}
% \include{clock-theory}
% \include{sv-ct-gw-dynamics}
%\include{environmental-models}

% %---
% \part{Pairwise Synchronization \\with Prior Localization}
\include{sync-parameter-estimation}
%\include{time-clock-control}
% \include{state-control}
\include{data-guided-sync}
%----
% \part{Ground \& Aerial Units:\\
% Simultaneous Localization and Synchronization}
% %---
\chapter{Kuiper Mobile Ground \& Aerial Networks}

In this chapter, we examine synchronization
for mobile units, augmenting the baseline
Kuiper network with ground and aerial vehicles.
%
Mobile units, in addition to the nominal SV/CT/GW
units, provide another set of challenges
as well as flexibility, to
devise novel synchronization protocols,
allowing for example, using SV in addition
to aerial vehicles to achieve better
synchronization at the CT during a GPS loss.
%---
Mobility can also be used to devise
effective differencing and diffusion-based
algorithms that further improve
synchronization performance across the
{\em augmented Kuiper} network; 
this is shown in Figure~\ref{mobility-1}.
\begin{figure}[ht]
\begin{center}
\includegraphics[width=.8\textwidth]{eps/clockeps/mobility.jpg}
\caption{Augmenting the Kuiper network with mobile aerial and ground nodes}\label{mobility-1}
\end{center}
\end{figure}

The synchronization problem with mobile units,
and potential use of differencing schemes for Kuiper, 
also allows us to consider the potential use of Kuiper for {\em navigation}.
%---
We provide the analysis on how to do this ``comm-for-nav'' could work for Kuiper after setting up the relevant relations and measurement schemes available from GPS and Kuiper. 
This is mainly done in the context of one-way time transfer;
extensions to two-way time transfer for mobile units
is also possible, left for future works.

Although mobile and aerial units have many uses in the
context of Kuiper network, this chapter is devoted
to providing the ground work for how mobility, and
the lack of position information for mobile units,
can be addressed in the one-way time transfer.
Other uses and algorithms for mobile ground and aerial Kuiper nodes
are left for future works.

\section{Observables for Synchronization}
%-----
The key construct in designing synchronization
protocols for any network, is what has commonly been
referred to as {\em observables} in the GNSS community.
Observables are measurements that are available
for each node to estimate the variable of 
interest, e.g., clock corrections.
%
For example, in the case of GPS, code
and carrier phase observables are the basic
constructs that are used for position and
clock offset estimation. This is shown
in Figure~\ref{gnss-receiver}.
\begin{figure}[ht]
\begin{center}
\scalebox{0.08}{\includegraphics{eps/clockeps/gnss-receiver-software-block.jpg}}\caption{GNSS receiver architecture} \label{gnss-receiver}
\end{center}
\end{figure}

Now, consider the analogous ``observables'' that
had been proposed by the Kuiper team for pairwise
SV-CT synchronization that is based on IEEE-1588.
\begin{figure}[ht]
\begin{center}
\scalebox{0.5}{\includegraphics{eps/clockeps/kuiper-syn-pair.jpg}}\caption{SV-CT sync by time stamping} \label{kuiper-team-sync}
\end{center}
\end{figure}

In this case, the observable is the time-stamping
of the SV signal with the purpose of comparing
\beas
\tau_{\mbox{\tiny SV}}(t) \quad \mbox{and} \quad \tau_{\mbox{\tiny CT}}(t);
\eeas 
the observable in this case is $\tau_{\mbox{\tiny SV}_i}(t_e)$,
that is received at $\mbox{CT}_j$ at time $t_r$, and
CT will aim to estimate $t_r$ using,
\beas
\tau_{\mbox{\tiny SV}_i}(t_e), \quad \tau_{\mbox{\tiny SV}_i}(t_e), \quad \mbox{and} \quad \widehat{\delta}_p,
\eeas 
where $\widehat{\delta}_p$ is the estimate of the propagation delay (more on this subsequently). In this case, the
CT clock offset is estimated by a sequence of measurements
(assuming that $\tau_{\mbox{\tiny SV}_i}(t)=t$ for all $i$ and
all $t$, and that have knowledge of the locations of
CT and SV precisely).
\begin{figure}[ht]
\begin{center}
\scalebox{0.15}{\includegraphics{eps/clockeps/sv-ct-sync.jpg}}\caption{SV-CT sync by time stamping} \label{sv-ct-sync}
\end{center}
\end{figure}
As we pointed out in the previous chapters,
if we use the ``dynamic'' model of the clock,
then Kalman filtering can reduce the number of
measurements for this scheme.

In a nutshell, if we can combine all our measurements
and assumptions (knowing the propagation delay),
we can establish an optimization problem 
that involves the unknown parameter (in this case, the clock 
bias $\alpha_j$ for $\mbox{CT}_j$), 
based on our model parameters and observations combined
in $\bm b$, i.e., 
\beas
\min_{\alpha_j} \|{\cal M}_e(\alpha_j, \bm b)\|
\eeas 
where $\bm b$ is our measurement/knowledge vector, $t_r$
in this case, is our unknown parameter;
the notation ${\cal M}_e$ is the error between what
the model ${\cal M}$ predicts and actual measurements;
this error is defined with respect to a particular norm (say, the 2-norm).
%
Notice that we can also have a batch processing setup,
where $t_r$ is only estimated after combining measurements
over many rounds.

Note that in the case of Kuiper SV-CT time synchronization,
we have assumed that we know the positions of CT and SV
(in a consistent coordinate frame). If we do not know
the position of the unit, then we had to formalize a 
different optimization model, such as
\beas
\min_{\alpha_j,{\bm x}_j} \| {\cal M}_e(\alpha_j, {\bm x}_j, \bm b)\|
\eeas 
where $\bm x_j$ would be the position of the unit (at a given time). In general, one can remove ambiguity in
estimating the underlying variables, if we have 
as many observations as the number of variables.
That is, if $\alpha_j$ and ${\bm x}_j$ are not
constrained, then we need four observations,
say, from four SVs. In order words, we can 
certainly use four SV timing signals to estimate
the clock bias {\em and} the position of the receiver.
This is shown in Figure~\ref{sv-mobile-sync}.

\begin{figure}[ht]
\begin{center}
\scalebox{0.1}{\includegraphics{eps/clockeps/moode-picture-mobile.jpg}}\caption{Estimating mobile unit position and clock offset using four range measurements from time-stamped signals from Kuiper SVs.} \label{sv-mobile-sync}
\end{center}
\end{figure}
%---------

If this time stamping (with respect to SV clocks) is
permissible, then it can be used as observables
to the base and the mobile units as well. In fact, we can
take this one step further and extend the base/mobile
or synchronization and positioning of a {\em network of mobile units}; these scenarios are shown in 
Figures~\ref{sv-base-mobile-sync}.
% and~\ref{sv-mobile-mobile-sync}.
%----
\begin{figure}[h]
\begin{center}
\scalebox{0.08}{\includegraphics{eps/clockeps/mood-base-mobile.jpg}}\;\;
\scalebox{0.061}{\includegraphics{eps/clockeps/mood-base-mobile-mobile.jpg}}
\caption{(Left) Estimating mobile unit position and clock offset using pseudorange measurements; (Right) estimating mobile unit position and clock offsets a network of relative observables} \label{sv-base-mobile-sync}
\end{center}
\end{figure}

% \begin{figure}[h]
% \begin{center}
% \scalebox{0.061}{\includegraphics{eps/clockeps/mood-base-mobile-mobile.jpg}}\caption{Estimating mobile unit position and clock offsets a network of relative observables} \label{sv-mobile-mobile-sync}
% \end{center}
% \end{figure}


\section{Mobile receiver position/timing architectures and assumptions}

In this section, we provide the basic architecture for
the mobile unit receiver position and timing algorithms.
The architecture allows inputs from GNSS when available
but more importantly, it will have a module to process
position and timing observables that can become
available from other Kuiper nodes, such as
SV time stamped signals. We also consider
how to use other mobile (ground or aerial)
units that can exchange position/timing
information with the given unit, as well as the so-called
signals of opportunity that might become available at a given
time. For example, in a recent work it has been
shown that carrier phase beat frequency can be
used for positioning of mobile units from 
LEO constellations such as Starlink--even
without a Starlink receiver. Similar analysis can
also be done for the Kuiper network.

\begin{figure}[ht]
\begin{center}
\scalebox{0.06}{\includegraphics{eps/clockeps/kuiper-mobile-1.jpg}}\caption{Mobile unit receiver architecture for positioning and time} \label{architecture-mobile}
\end{center}
\end{figure}

Let us provide distinct algorithms for position and timing
for mobile units that are suitable under different assumptions;
we generically call these algorithms PNT for mobile units.

\subsection{Mobile PNT with GNSS}

The baseline assumption in this case is that the mobile unit has
a GNSS/GPS receiver, providing timing corrections and
pseudorange information to the mobile unit.

\begin{enumerate}
\item GNSS pseudoranges (from independent directions) can be processed by the mobile unit receiver to output PNT information for a stationary case; these measurements can also be updated for a slowly moving unit.
\item If the mobile unit is moving, then we can integrate GNSS measurements with an inertial navigation system (INS). The INS uses
the information from the inertial measurement units (IMUs) {\em in addition to} the pseudorange measurements from GNSS to update the unit position as well as correct for its clock offset and skew. The setup consists of a nonlinear (extended) Kalman filter built on the mobile unit dynamics (position/velocity and orientation when applicable) as well as the clock dynamics (for example, the second order model that we had discussed previously). In this case, we can also update the clock noise terms depending on the environmental conditions (e.g., outside temperature).
\item We will discuss the notion of geometric dilution of precision (GDOP) shortly, that can be computed for the globe. If GDOP is poor for certain locations on earth, the GNSS/INS integration can also prove to be useful for mobile units that can operate with intermittent GPS connectivity.
%----
\item In the case when GDOP is poor from the GNSS signal,
one can also consider using other Kuiper nodes (SV, base station, aerial and other ground units) to further improve 
the nominal GDOP. We understand the basic geometry
that dictates GDOP for GNSS that can also
be used to reason as to when this augmentation makes sense
and useful. The background on such schemes
are (double/triple) differencing as well as diffusion-based
protocol that can further improve the PNT accuracy and precision
for the mobile unit.
\end{enumerate}

\subsection{Mobile PNT with Kuiper with partial or no GPS} 

There has been a lot of discussions in space-based
PNT and communication on using commercial space assets
for improving, augmenting, or even replacing GPS.
This point of view is particularly relevant when GPS
signals becomes unavailable or disrupted. Although,
the baseline Kuiper mode utilizes GPS, there are compelling
reasons for examining a ``backup mode'' for Kuiper,
one that only relies on partial or no GPS/GNSS.

In this chapter, we examine how Kuiper time-stamping
schemes can be used in the context of positioning
and timing. There are many variations to this setup;
hence, our primary objective is to show that
in principle, Kuiper can provide timing and positioning
functionality, augmenting GPS, or using partial or no
GPS; this is shown both for ground rovers as well
as aerial units that we refer to as Kuiper aerial vehicles
(KAVs) or drones.

\begin{enumerate}
\item We derive how the time-stamping scheme between
Kuiper SV and the mobile units can mimic pseudorange
measurements obtained from GNSS. We can even augment
time stamping with carrier beat frequency to
create ``two'' loops for pseudorange estimation that
further improve PNT accuracy and precision.
\item Kuiper SV time stamping can also be replicated
for aerial units; we should some simulations for 
this setup in this chapter.
%
\item Similar to the GPS case, we can augment time stamping
from Kuiper SV with INS to improve the time-stamping schemes
and further improve PNT performance at the mobile unit.
This augmented Kuiper/IMU can be effective for
ground and aerial vehicles that do not have access to 
sufficient satellite time-stamped signals at all times.
%
\item Kuiper time stamping can also be embedded in (double/triple) differencing with other units in the Kuiper network (ground or aerial vehicles) as well as with diffusion-based schemes to further improve the accuracy and precision of mobile units' PNT.
\end{enumerate}

% \subsection{Mobile PNT with GNSS and Kuiper}

% This can be an attractive option in certain regions
% of the globe- for example, Kuiper SV geometry can
% complement the poor coverage by GNSS at certain locations.
% more on this ...

\subsection{Mobile PNT with Signals of Opportunity}

This was suggested during our weekly discussion 
and might of interest to examine further.
By signals of opportunity, we means signals
that had not been intended specifically for PNT,
but can nevertheless be used for such purposes.
For example, if the mobile unit can recognize 
a landmark and could estimate its
relative position with respect to that landmark, then
one can examine to what extend this ``signal of opportunity''
can be used either for position estimation, or
where to look in the sky for a Kuiper SV.
Another example, is to use the carrier beat frequency
at the receiver to estimate ranges- that when independent,
can be used to position estimation.



\begin{figure}[h]
\begin{center}
\scalebox{0.08}{\includegraphics{eps/clockeps/kuiper-mobile-2.jpg}} 
\caption{Position estimation by monitoring the signal strength from one SV with fixed $\delta t_e$ assuming the SV antenna is not articulated.} \label{SV-signal-strength} 
\end{center} 
\end{figure}


\begin{enumerate}
\item Suppose the only measurement available at the mobile unit is the received signal power from a given Kuiper SV. It is then conceivable that the receiver knows at least that it
is in the particular satellite footprint. If time stamping
is also embedded in the signal, then we can use
the time information to know the position of the SV and
a one-dimensional pseudorange with respect to the SV.

\item If the SV-mobile unit comm is via a spot beam,
then the coverage distance is more refined, and thus
more useful for positioning. For example, for a SV in 
630 km orbit, with a one degree half-beam angle,
the coverage area is around 380 square km (radius 
from the sub satellite point is about 11 km).

\item In order to further refine the above uncertainty,
we can now look a sequence of footprints from on satellite. 
For example, if the signal strength drops drastically over 
a time interval of length $\delta t_e$, then the receiver 
can estimate its
position to be in the overlap region between the footprint
just before the signal attenuation; this is shown in 
Figure~\ref{SV-signal-strength}. For a SV in 630 km circular
orbit with a half-beam of 1 degree, the width of the overlap is
\beas
(2\pi/\mbox{period}) \times R_e \times \delta t_e
\eeas
so if $\delta t_e$ is about a second, then the width of the overlap
is about 6 km.

\item One can also be more creative with using the multi-beam
geometry; discuss the multi-beam geometry from the single
SV from the patent found online (search ``geolocation communication'').

\end{enumerate}

In the following, we consider the case-by-case scenarios
for mobile unit PNT; along the way, we also
address the following questions:
\begin{shaded*}
\begin{enumerate}
\item What are the parallels between GPS receiver processing and its potential applications for Kuiper? 
For example, can you use carrier phase information in addition to time-stamping to further improve the synchronization performance?
\item How effective can differencing schemes be for augmented Kuiper nodes? In particular, we will examine
how differencing schemes between the base station and mobile nodes, can be used to improve synchronization
performance across the network?
\item To what extend can the Kuiper signals be used for navigation? This issue has been examined for SpaceX signal in recent paper (using carrier phase differencing) and extended Kalman filtering.
\end{enumerate}
\end{shaded*}

% \section{Mobile Unit PNT: case by case scenarios}

% In this section, we delve into various assumptions
% that we alluded to above, before propose distinct
% algorithms that can be used for mobile unit PNT.

\subsection{Positioning based on the received signal strength}

In the next section, we examine the situation when
time stamped information can be relayed to the
mobile units by say the Kuiper SV. Such measurements
can then used to define observables that subsequently
can be solved for position and timing. The more
basic question that we like to examine in this
section is the extent by which one can use received
signal strength from a Kuiper SV for position,
without any time stamping information.

Suppose that the rover has established a communication link
with a Kuiper SV.\footnote{Tim asked what happens if the rover has no information about its position and can not determine where in the sky to look for the SV. There are a few approaches to this scenario that might be worth examining further. Of course, if the rover had some knowledge about its position at {\em some} time in the past, then it could use its inertial navigation system (INS) to update its position information using onboard inertial measurement units (IMUs) provided periodic updates from GPS or Kuiper units. Assuming that this information/update mechanism is not available to the rover, and that it can not receive this information (or its noisy version) from other units- then we can use onboard cameras to estimate the rover's position based on some landmarks or positions of sun or moon in the sky- for example, performing simultaneous localization and mapping using onboard cameras and IMUs. The other approach is have a scanning mechanism for the rover antenna; we can use the Kuiper geometry to devise an optimized scanning profile that increases the probability of establishing a link with a Kuiper SV.}
The first question is what additional information does this
link establishment provide for the position/timing estimation.
To encode this information in the optimization problem,
consider the geometry shown in Figure~\ref{link-geometry}.
\begin{figure}[ht]
\begin{center}
\scalebox{0.15}{\includegraphics{eps/clockeps/sat-spot-beam-geometry.jpg}}  
\caption{Beam angle and coverage area geometry} \label{link-geometry}
\end{center}
\end{figure}
%------
With reference to Figure~\ref{link-geometry},
the coverage area is the function of the angle
$\gamma$, namely,
\beas
A_{\mbox{\tiny cov}} = 2 \pi R_{\mbox{\tiny e}}^{2} (1 - \cos \gamma)
\eeas 
that is related to the beam (half) angle $\alpha$.
If we start with $\alpha$, then we can use the geometry
to obtain $\gamma$ and the coverage area. To do this, from
the law of cosines, we observe that,
\beas
R_s^{2} = R_{\mbox{\tiny e}}^2 + r_{\mbox{\tiny SV}}^2 - 2 r_{\mbox{\tiny SV}} R_{\mbox{\tiny e}} \cos \gamma
\eeas 
where $R_{\mbox{\tiny e}}$ is the radius of Earth
and $r_{\mbox{\tiny SV}}$ is the distance from the
center of earth to SV (that is, the radius of Earth added to the SV altitude). In the meantime, we can also write,
\beas
R_{\mbox{\tiny e}}^2 = R_{s}^2 + r_{\mbox{\tiny SV}}^2 - 2 r_{\mbox{\tiny SV}} R_{s} \cos \alpha;
\eeas 
given $\alpha$ we can thus find $R_s$ from the last
equation, and then use the previous relationship
to find $\gamma$ and subsequently the coverage area.
For Kuiper the coverage area would be about $1500$ square
kilometers using the half-angle beam of 2 degrees.
% \beas
% \alpha = \sin^{-1} R_e/r
% \eeas 
So this would be useful if the initial GDOP
uncertainty provides an uncertainty that is
larger than this characterization.

In order to provide an algorithm for how
``side'' information can be used for estimating
the position of the rover, let
${\widehat {\bm r}}_{\mbox{\tiny SV}}$ and ${\widehat {\bm r}}_{\mbox{\tiny R}}$ denote, respectively the normalized vectors
to the Kuiper SV and the rover; we have to fold in the
fact that rover's position vector will be expressed in
ECEF frame (rotating with Earth) whereas SV position
vector is generally expressed in ECI (inertially fixed).
So there is a (time varying coordinate transformation),
say, $U(t)^\top$ that puts ${\widehat {\bm r}}_{\mbox{\tiny SV}}$ in
the ECEF frame at time $t$. Using this transformation,
we then expressed that the rover is in the SV coverage
zone by the quadratic inequality,
\beas
{\widehat {\bm r}}_{\mbox{\tiny SV}}(t)^\top \, U(t)^\top \, {\widehat {\bm r}}_{\mbox{\tiny R}}(t) \leq \cos \gamma
\eeas
where $\gamma$ is dictated by the beam angle $\alpha$
as discussed previously; this is 
shown in Figure~\ref{cone-constraint-visibility}.
\begin{figure}[ht]
\begin{center}
\scalebox{0.15}{\includegraphics{eps/clockeps/sat-spot-beam-geometry2.jpg}} 
\caption{Position inequality induced by link existence} \label{cone-constraint-visibility}
\end{center}
\end{figure}

Here is now the fundamental observation- if have
a communication link over a certain time interval, just
knowing that there is a link over a given interval,
we can find the feasible point for
\beas
{\widehat {\bm r}}_{\mbox{\tiny SV}}(t_1)^\top \, U(t_1)^\top \, {\widehat {\bm r}}_{\mbox{\tiny R}} \leq \cos \gamma\\
{\widehat {\bm r}}_{\mbox{\tiny SV}}(t_2)^\top \, U(t_2)^\top \, {\widehat {\bm r}}_{\mbox{\tiny R}} \leq \cos \gamma\\
{\widehat {\bm r}}_{\mbox{\tiny SV}}(t_3)^\top \, U(t_3)^\top \, {\widehat {\bm r}}_{\mbox{\tiny R}} \leq \cos \gamma\\
\!\!\!\!\vdots\\
{\widehat {\bm r}}_{\mbox{\tiny SV}}(t_n)^\top \, U(t_n)^\top \, {\widehat {\bm r}}_{\mbox{\tiny R}} \leq \cos \gamma
\eeas
assuming that the rover has been stationary during
this interval (we have removed the time dependency
for ${\bm r}_{\mbox{\tiny R}}$).

In fact, if we add this set of inequalities to our
original estimation problem, we obtain a nice
optimization of the form,
\beas
 \min_{{\bm x_{\mbox{\tiny R}}}}} \, \|\widehat{\rho}_{1:4}- \widehat{\rho} ({\bm x_{\mbox{\tiny R}}}})\| \\
{\widehat {\bm r}}_{\mbox{\tiny SV}}(t_i)^\top \, U(t_i)^\top \, {\widehat {\bm r}}_{\mbox{\tiny R}} \leq \cos \gamma, \, \quad i=1,2,\ldots, n,
\eeas 
after $n$ rounds of communication; this is depicted in 
Figure~\ref{sequence}.
%---
% \\~\\
% \noindent{\color{blue} The code for solving this optimization problem.}
% \\
\begin{figure}[ht]
\begin{center}
\scalebox{0.2}{\includegraphics{eps/clockeps/sat-spot-beam-geometry3.jpg}} 
\caption{Using a sequence of beam images to estimate position of the rover} \label{sequence}
\end{center}
\end{figure}

Now, another aspect of this geometry is the scenario
when we have multiple beams on the SV (a search online
with geolocation via satellite communication pointed
to a patent application- review with the team!)
%---
Then the overlap of the corresponding cones on the ground,
even at a single time, can provide a set of
quadratic inequalities that augmented,
the time stamping, for position estimation of the rover.

% More fundamentally, this chapter examines various
% means of estimating position and timing for
% mobile ground and aerial units.

% The overall architecture that we will consider is
% depicted in 


% The idea is to what extend can you determine
% the position of the unit just by knowing when 
% the satellite is visible.


% it is the rate by 

% \beas
% dA/dt = \frac{dA}{d\delta} \, \frac{d \delta}{dt}
% \eeas 


%--------
\section{Position/clock bias estimation for (slow) mobile units using time stamping}
In this section, we examine how we can estimate the
position and clock bias of a mobile unit using 
SV time-stamped measurements. We then examine the
corresponding uncertainties depending on how these
measurements have come about, e.g., from 
$n$ number of SVs at a specific time or say, one SV over 
$n$ subsequent times.

In order to distinguish between time-stamping based
positioning and timing and those based on say, code/carrier phase
tracking using in GPS (or carrier phase for Kuiper), 
we first define the notion of (true)
TS (time-stamping) observable (between the space vehicles SV and the mobile unit, or the rover, R; this
observable which we call the ideal pseudorange
assumes the form,
\beas
\rho(t_{r_1},t_{e_1}}) = c \, \{ t_{r_1} -t_{e_1} + \delta t_{\mbox{\tiny} R} \},
\eeas 
where $\delta t_{\mbox{\tiny} R}$ is the rover clock bias
at $t_{r_1}$ (for the moment we assume the corresponding
SV clock bias is negligible or has been accounted for);
$c$ denotes the speed of light.


\begin{figure}[ht]
\begin{center}
\scalebox{0.15}{\includegraphics{eps/clockeps/moode-picture-mobile-singleSV.jpg}} 
\caption{Fundamental limitation position and timing using single SV measurements} \label{PNT-singleSV}
\end{center}
\end{figure}


The actual TS-pseudorange measurement however would be
\beas
\widehat{\rho}(t_{r_1},t_{e_1}}) = c \, \{ \tau_{\mbox{\tiny R}}(t_{r_1})} - \tau_{\mbox{\tiny SV}}(t_{e_1})} \}
\eeas 
where essentially the SV and rover clocks are used to measure
the difference between the emission and reception times
of the time-stamped signal from SV to the rover.
%
The pseudorange expression can now be expanded as,
\beas
\widehat{\rho}(t_{r_1},t_{e_1}}) = \rho(t_{r_1},t_{e_1}}) + \rho_{\mbox{\tiny trop}} + \rho_{\mbox{\tiny ion}} + \rho_{\mbox{\tiny MP}} + \rho_{\mbox{\tiny random}}
\eeas 
where
\beas
{\rho}(t_{r_1},t_{e_1}}) = \|\bm r_{\mbox{\tiny SV}}(t_{e_1}) - \bm r_{\mbox{\tiny R}}(t_{r_1}) \| + c \, \delta t_{\mbox{\tiny R}}(t_{r_1}),
\eeas 
with vectors expressed with respect to a consistent
coordinate frame, and $\rho_{\mbox{\tiny trop}}$, $\rho_{\mbox{\tiny ion}}$, $\rho_{\mbox{\tiny MP}}$, $\rho_{\mbox{\tiny random}},$ are the corresponding (length-scale) delays due to
troposphere, ionosphere, multipath, and random effects.

Now let 
\beas
\widehat{\rho}_{1:4}= \leftm{c} \widehat{\rho}(t_{r_1},t_{e_1}}) \\
				\widehat{\rho}(t_{r_2},t_{e_2}}) \\
				\widehat{\rho}(t_{r_3},t_{e_3}}) \\
				\widehat{\rho}(t_{r_3},t_{e_4}}) \rightm;
\eeas 
assume that we know the estimate of $r_{\mbox{\tiny SV}}(t_{e_1})$ as
\beas
\widehat{\bm r}_{\mbox{\tiny SV}}(t_{e_1}) = \bm r_{\mbox{\tiny SV}}(t_{e_1}) + \bm e_{\mbox{\tiny E}}
\eeas 
where $e_{\mbox{\tiny E}}$ denotes the ephemeris
estimation error.
Our estimation problem is now estimating 
$\widehat{\bm r}_{\mbox{\tiny R}}(t)$ for
some $t \in [t_{e_1},t_{r_4}]$; this will be pursued
by minimizing problem,
\beas
\min_{\widehat{\bm r}_{\mbox{\tiny R}}(t)} \|
\leftm{c}
\widehat{\rho}(\bm r_{\mbox{\tiny R}}(t_{r_1}))\\
\widehat{\rho}(\bm r_{\mbox{\tiny R}}(t_{r_2}))\\
\widehat{\rho}(\bm r_{\mbox{\tiny R}}(t_{r_3}))\\
\widehat{\rho}(\bm r_{\mbox{\tiny R}}(t_{r_4})) \rightm
- \widehat{\rho}_{1:4} \|
\eeas 
with respect to $\bm r_{\mbox{\tiny R}}$; noting that
$\widehat{\rho}$ is a nonlinear function of 
$\bm r^{1}_{\mbox{\tiny R}}$, $\bm r^{2}_{\mbox{\tiny R}}$,
$\bm r^{3}_{\mbox{\tiny R}}$ at the respective times.

Then the problem at hand looks like,
\beas
\min_{{\bm x_{\mbox{\tiny R}}}}} \, \|\widehat{\rho}_{1:4}- \widehat{\rho} ({\bm x_{\mbox{\tiny R}}}})\|
\eeas 
where
\beas
\bm x_{\mbox{\tiny R}} = \leftm{c} {\bm r_{\mbox{\tiny R}}}}  \\ c\, \delta t_{\mbox{\tiny R}} \rightm 
\eeas
and 
\beas
\widehat{\rho}(\bm x_{\mbox{\tiny R}}) = \leftm{c} \|\bm r_{\mbox{\tiny SV}}(t_{e_1}) - \bm r_{\mbox{\tiny R}}(t) \| + c \, \delta t_{\mbox{\tiny R}}(t) \\\|\bm r_{\mbox{\tiny SV}}(t_{e_2}) - \bm r_{\mbox{\tiny R}}(t) \| + c \, \delta t_{\mbox{\tiny R}}(t) \\ \|\bm r_{\mbox{\tiny SV}}(t_{e_3}) - \bm r_{\mbox{\tiny R}}(t) \| + c \, \delta t_{\mbox{\tiny R}}(t) \\ \|\bm r_{\mbox{\tiny SV}}(t_{e_4}) - \bm r_{\mbox{\tiny R}}(t) \| + c \, \delta t_{\mbox{\tiny R}}(t) \rightm + \mbox{\bf vec} \, (\rho_{\mbox{\tiny trop}} + \rho_{\mbox{\tiny ion}} + \rho_{\mbox{\tiny MP}} + \rho_{\mbox{\tiny random}})
\eeas 
where $t$ is between $t_{e_1}$ and $t_{r_4}$ with the 
implicit assumption that the mobile unit has not
changed its position in this interval. If it has,
then we can use the accelerometer to to relate
$\bm r_{\mbox{\tiny R}}(t)$ from one time stance to
another.
%---

Note that if we have two or more rovers (say $n)$, or aerial vehicles (say $m$), then can use the designation 
the position and their respective clock parameters as,
\beas
\bm x_{\mbox{\tiny R}_1} = \leftm{c} {\bm r_{\mbox{\tiny R}_1}}}  \\ c\, \delta t_{\mbox{\tiny R}_1} \rightm, \, \quad \bm x_{\mbox{\tiny R}_2} = \leftm{c} {\bm r_{\mbox{\tiny R}_2}}}  \\ c\, \delta t_{\mbox{\tiny R}_2} \rightm, \, \ldots, \, \leftm{c} {\bm r_{\mbox{\tiny R}_n}}}  \\ c\, \delta t_{\mbox{\tiny R}_n}\rightm,
\eeas
and
\beas
\bm x_{\mbox{\tiny A}_1} = \leftm{c} {\bm r_{\mbox{\tiny A}_1}}}  \\ c\, \delta t_{\mbox{\tiny A}_1} \rightm, \, \quad \bm x_{\mbox{\tiny A}_2} = \leftm{c} {\bm r_{\mbox{\tiny A}_2}}}  \\ c\, \delta t_{\mbox{\tiny A}_2} \rightm, \, \ldots, \, \leftm{c} {\bm r_{\mbox{\tiny A}_m}}}  \\ c\, \delta t_{\mbox{\tiny A}_m}\rightm;
\eeas
Then depending on the information transfer between
the nodes, we can not only use the information about 
$\widehat{\rho}(\bm x_{\mbox{\tiny R}_1})$
and $\widehat{\rho}(\bm x_{\mbox{\tiny R}_2})$,
for their individual state estimation, but 
potentially use the inter-rover (or rover-aerial)
timing and distance information for estimation. 
This would also be the case with the base
station, that is often assumed to be at a {\em known} location.

For now, let us consider the single rover case;
the general approach for solving the position
and timing (of the rover) proceeds 
as follows: Let us have an estimate of the position and clock offset for the rover, namely $\bm x_{\mbox{\tiny R}}(0)$;
then we update these estimates using the recursion,
\begin{shaded*}
\beas
\bm x_{\mbox{\tiny R}} (k+1) = \bm x_{\mbox{\tiny R}}(k) + (H(k)^{\top} H(k))^{-1} H(k) \, (\widehat{\rho}_{1:4} - \widehat{\rho}(\bm x_{\mbox{\tiny R}} (k))), \; \mbox{for} \; k=0, 1, 2, \ldots
\eeas 
\end{shaded*}
where,
\beas
H(k) &=& \nabla \widehat{\rho}(\bm x(k)) \\
&=& \leftm{cccc} 
(\bm r^{1}_{\mbox{\tiny SV}}(t_{e_1}) - \bm r^{1}_{\mbox{\tiny R}}(k))/p_1 & (\bm r^{2}_{\mbox{\tiny SV}}(t_{e_1}) - \bm r^{2}_{\mbox{\tiny R}}(k))/p_1 & (\bm r^{3}_{\mbox{\tiny SV}}(t_{e_1}) - \bm r^{3}_{\mbox{\tiny R}}(k))/p_1& 1\\[.05in]
(\bm r^{1}_{\mbox{\tiny SV}}(t_{e_2}) - \bm r^{2}_{\mbox{\tiny R}}(k))/p_2 & (\bm r^{2}_{\mbox{\tiny SV}}(t_{e_2}) - \bm r^{2}_{\mbox{\tiny R}}(k))/p_2 & (\bm r^{3}_{\mbox{\tiny SV}}(t_{e_2}) - \bm r^{3}_{\mbox{\tiny R}}(k))/p_2 & 1\\[.05in]
(\bm r^{1}_{\mbox{\tiny SV}}(t_{e_3}) - \bm r^{1}_{\mbox{\tiny R}}(k))/p_3 & (\bm r^{2}_{\mbox{\tiny SV}}(t_{e_3}) - \bm r^{3}_{\mbox{\tiny R}}(k))/p_3& (\bm r^{3}_{\mbox{\tiny SV}}(t_{e_3}) - \bm r^{3}_{\mbox{\tiny R}}(k))/p_3& 1\\[.05in]
(\bm r^{1}_{\mbox{\tiny SV}}(t_{e_4}) - \bm r^{1}_{\mbox{\tiny R}}(k))/p_4 & (\bm r^{2}_{\mbox{\tiny SV}}(t_{e_4}) - \bm r^{2}_{\mbox{\tiny R}}(k))/p_4 & (\bm r^{3}_{\mbox{\tiny SV}}(t_{e_4}) - \bm r^{1}_{\mbox{\tiny R}}(k))/p_4 &  1
\rightm 
\eeas
and
\beas
p_i= \|\bm r_{\mbox{\tiny SV}}(t_{e_i}) - \bm r_{\mbox{\tiny R}}(k)\|, \quad i = 1, 2, 3, 4.
\eeas
In the GPS case, we generally need at least four measurements to
resolve for the position {\em and} receiver clock offset
but more measurements would further improve the quality of these
estimates-- note that in the above expression
for $H(k)$, our notation suggests that we are using
four measurements but in general we can use
more measurements when available.

The key observation is that the rows of the matrix
$H(k)$ (that are obtained from the linearization
of pseudorange equation) are comprised of the normalized vectors from the mobile unit to the SVs (or the base station), 
augmented by 1's, that pertain to the clock correction,
Note that these vectors are generally represented in the earth-centered/earth fixed frame
(ECEF) so appropriate coordinate transformation should be
used. Hence, in a nutshell, the rows of the matrix
of interest relate to the independence structure of
directions from the mobile unit to the SVs; see~Figure~\ref{sv-sync}.
\begin{figure}[ht]
\begin{center}
\scalebox{0.18}{\includegraphics{eps/clockeps/moode-picture-mobile-geometry}} 
\caption{Geometry of SV/Mobile unit synchronization} \label{sv-sync}
\end{center}
\end{figure}

This iteration usually converges after a few iterations
if the matrix $H$ (for various values of $k$) is well-conditioned.
Once the iterates converge, we have
\beas
(H^{\top} H)^{-1} H \, (\widehat{\rho}_{1:4} - \widehat{\rho}(\bm x_{\mbox{\tiny R}})) =0
\eeas
where the matrix $H$ is the evaluation of $H(k)$ at 
$\bm x_{\mbox{\tiny R}}$.

If the measure and delay estimates have covariance $\Sigma$,
the the covariance of the error estimate (in position and timing)
is,
\beas
(H^{\top} H)^{-1} H^\top \Sigma H (H^\top H)^{-1}
\eeas
and in particular, when $\Sigma = \sigma^{2} I$, then
the delay/measurement error covariance are amplified as 
$D \sigma^2$ where,
\beas
G= (H^{\top} H)^{-1} .
\eeas
The matrix $G$ is often called the {\em Geometric Dilution of Precision} or GDOP; this matrix is essentially a 
multidimensional sensitivity measure of how variations
in the measurements effect the variation of the estimates
(in position and time). One can even go one step
further and consider,\footnote{The subscripts refer to the entries of the matrix.}
\beas
\sqrt{G_{33}}, \sqrt{G_{11}+D_{22}}, \; \mbox{and} \; \sqrt{G_{11}+G_{22}+G_{33}},
\eeas 
that signify, respectively, the sensitivity of the estimates in altitude, planar, and total position. This is shown 
geometry below- GDOP captures the uncertainty 
in distinct directions due to the estimation setup;
this is shown in Figure~\ref{GDOP-geometry}.
%
\begin{figure}[ht]
\begin{center}
\scalebox{0.15}{\includegraphics{eps/clockeps/GDOP}}
\caption{Geometry of the dilution of precision}
\label{GDOP-geometry}
\end{center}
\end{figure}

Note that the matrix $G$ is a function of position of the receiver with respect to the satellites; 
the conditioning $G$
is at the heart of the issue--so we can evaluate
$H$ at $\bm x_{\mbox{\tiny R}}$ with respect to
say four values for $\bm r_{\mbox{\tiny SV}}$'s.
%--- 
Now comes the fundamental issue with estimating position and
timing based on four $\bm r_{\mbox{\tiny SV}}$ 
for a single SV, say at four distinct time stamps.
The issue is that single SV orbits are planar in the
earth-fixed inertial frame. But this then means that
no matter how many time stamps we have from a single satellite,
we will have the linear dependence between these vectors
that does not facilitate position and timing estimation
from single satellite measurements. As such, depending
on the direction where the position estimation is
necessary (say "north" or altitude), one might not
be able to resolve position and timing, as such,
we should rely on an ``augmentation'' of the SV measurements
with other measurements, says from a base station or
a mobile unit; this is depicted in Figure~\ref{sv-sync-mobile}.
%---
\begin{figure}[ht]
\begin{center}
\scalebox{0.15}{\includegraphics{eps/clockeps/moode-picture-mobile-geometry-2}}
\caption{Geometry of mobile ground and aerial units positioning and time synchronization}
\label{sv-sync-mobile}
\end{center}
\end{figure}
%\vspace{.1in}
% \\
% \noindent {\bf Example:} To demonstrate the
% setup above, consider a Kuiper SV sending
% time stamped signals to the receiver--
% this has been done in Jupyter Notebook and
% the result will be included here.


% \subsection{Another variation on position and time estimation}

% In the previous section we explored how the time stamping
% information can be used for range estimates (pseudorange)
% and how a collection of independent range estimates
% leads to least squares problem of the form
% \beas
% \min_{{\bm x_{\mbox{\tiny R}}}}} \, \|\widehat{\rho}_{1:4}- \widehat{\rho} ({\bm x_{\mbox{\tiny R}}}})\|
% \eeas 
% where ${\bm x_{\mbox{\tiny R}}}$ contains the 
% position and the clock information for the ``rover,''
% e.g., mobile unit. It is the quality of this estimate
% that puts a fundamental bound on the estimation accuracy-
% and one can adopt augmenting pseudorange observables
% with other measurements, similar to differential GPS or
% information for a ground base station or an aerial vehicles.
% (this ``network'' approach will be examined in more details).
% However, for LEO satellites and in particular Kuiper, there
% is another type of information that can be useful for
% position estimation.


% \subsubsection{Block diagrams}

% \noindent {\color{blue} Block diagram for the Kuiper positioning and timing for mobile and aerial units.}\\


% \noindent {\color{blue} Block diagram for the Kuiper/GPS positioning and timing for mobile and aerial units.}\\


% \noindent {\color{blue} Compute of the GDoP for Kuiper alone,
% Kuiper+GPS at ever point on the globe; identify the variations}\\\\
% \noindent {\color{blue} another issue that we need to consider
% is when clocks have a skew; if you include them 
% in the derivation of the pseudorange, then
% they show up in the pseudorange equation; in this
% case, it seems that we can use a two way communication 
% to obtain the skew first, then proceed with 
% setting up the estimation as a bi-level problem:
% estimate the mobile unit clock skew, followed by position
% and timing- a Kalman filter will prove instrumental to
% streamline this process.} \\~\\

% % If we assume that the receiver clock offset has a 
% % frequency bias, then
% % \beas
% % \delta_{\mbox{\tiny R}}(t_{r_2}) = \delta_{\mbox{\tiny R}}(t_{r_1}) + \beta (t_{r_2} - t_{r_1})
% % \eeas 
% % with $t_{r_2}-t_{r_1} \approx t_{e_2}-t_{e_1}$
% % leads to addition of one more parameter to the optimization problem.\\~\\

% % Let us now consider an example to underscore the
% % importance of the Kuiper geometry, or the 
% % use of a base or aerial platform in the estimation
% % process. More sophisticiated estimation/filtering
% % algorithms can be adopted of course, but the setup
% % above can be used to highlight the fundamental issues
% % in position and clock offset computaiton (and correction).\\~\\

% % if we have clock skew as well we can do a
% % two way time stamping if needed.

% % \noindent {\bf Example:} see the Jupyter notebook for the estimation problem without the Kalman filter and
% % with Kalman filter\\~\\

% % % First, let us provide the key algorithms steps
% % % for GPS, differential GPS, as well as augmented
% % % GPS-INS (Inertial Navigation Systems) for clock offset
% % % estimates and position estimation for mobile units.


% % altitude known and unknown\\~\\


% \noindent {\bf Example:} how does this compare with GPS


% \subsubsection{GDOP heat map for Kuiper over the globe}

% given a particular visibility cone, we can compute,
% for various phases of Kuiper constellation, 
% the dilution of precision for every land mass across
% the globe in terms of the heat map.

\section{Overarching theme for positioning and timing networks}

A network consists of nodes that have knowledge on their
positioning and timing (or that of other nodes) 
with respect to a global reference frame.
%----
Each node can share information (absolute or relative) timing/position
information with other nodes but this information is noisy and delayed.
The various algorithms proposed operate at the node level: given
each node uncertainty (or lack of knowledge) of its own timing and position,
aim to reduce the uncertainty of timing/position at various nodes-
and in the timing case, synchronize. These nodes can be Kuiper SVs,
CTs, GWs, as well as aerial and mobile units. In the case of IoT
devices, energy constraints can also be folded in.

\section{Standalone Kuiper Positioning and Timing}

In this section, we consider standalone Kuiper PNT
for mobile units. The basic building block for
this scheme is time-stamped observables available
to the unit from Kuiper SVs. 
Moreover, time-stamped observables can
potentially be augmented with an inertial navigation system 
(INS) including measurements available from the inertial measurement units (IMUs).

\subsection{CT vs. mobile unit clock synchronization performance}
Prior to developing position and timing algorithms for
mobile units (which we subsequently refer to as rovers),
utilizing the dynamics of the vehicle
and its clock, let us examine how to assess the performance
of synchronization algorithms {\em in the absence} of such algorithms.
In particular, we examine the inherent synchronization
performance when we decouple time and position estimation
for the rover. 

Suppose that the rover receives a timing signal
from a Kuiper SV (tagged as $\tau_t_{e}$) and the SV id. 
To streamline the algebra for now, we assume that
$\tau_{\mbox{\tiny SV}}(t_{e}) \approx t_{e}$
that is the SV clock is accurate.
%
Then the SV id and the time stamp can be used to 
know $\bm r_{\mbox{\tiny SV}}(t_{e})$; this
position vector can be represented in the ECI
or ECEF frames.

If the rover's location was fixed on earth's surface,
then receiving this time tagged signal at $t_{r}$,
measured as $\tau_{\mbox{\tiny R}}(t_{r})$, would then
allow the rover to formulate an observation equation
for the propagation time, i.e,
\bea
\widehat{t}_p \approx \| \bm r_{\mbox{\tiny SV}}(t)- \bm r_{\mbox{\tiny R}}(t) \|/c ,
\eea 
where $\widehat{t}_p$ is the estimate of
the signal propagation time and $c$ is the signal
propagation speed; note the ambiguity in $t$,
as the rover clock does not have access to the universal time.
The signal propagation time is then used to calibrate the clocks. Any error in the signal propagation is then reflected in the 
clock error.
%
Hence, even at this point, for every meter position
uncertainty (even if the clocks are accurate)
we incur about 3.33 ns of time uncertainty. 

When the clocks are {\em not} accurate, the situation
becomes more interesting, as the emission and
receptions times, and the corresponding position
vectors have to all be adjusted and compensated for.
That is, suppose that $t_{e}$ is accurate but
$t_r$ is measured by the rover clock at $\tau(t_r)$
and we like to use,
\bea
\widehat{t}_p \approx \| \bm r_{\mbox{\tiny SV}}(t_e)- \bm r_{\mbox{\tiny R}}(t_r) \|/c ,
\eea
Then assuming an algebraic model for the clock onboard
the rover,
\beas
\tau_{\mbox{\tiny R}}(t_r) = \alpha + \beta t_r 
\eeas 
we can see that the timing error would be related
to the difference,
\beas
\|\bm r_{\mbox{\tiny R}}(t_r) - \bm r_{\mbox{\tiny R}}(\tau_{\mbox{\tiny R}}(t_r))\| \approx \| \bm v_{\mbox{\tiny R}} (t_r) (t_r- \tau_{\mbox{\tiny R}}(t_r))\|;
\eeas 
so even in this simple setting, the estimate of
propagation time needed for clock correction, 
are related to both the translational and clock dynamics
of the rover (the bias and skew).
\begin{figure}[ht]
\begin{center}
\scalebox{0.05}{\includegraphics{eps/clockeps/rover-position-error}}
\caption{Kuiper mobile unit position/timing using SV time stamping with position uncertainty}
\label{time-stamping-position-uncetatinty}
\end{center}
\end{figure}
For example, if we only consider range ambiguity, 
from Figure~\re{time-stamping-position-uncetatinty},
and using the law of cosines, it follows that,
\beas
d_3^2= d_1^2 + d_2^3 - 2 d_1 d_2 \cos \theta
\eeas 
where $d_2 = v \Delta t$. 
Hence, the percent error in the $t_p$ computation is 
\beas
\frac{d_3}{d_1} = \sqrt{1+ \frac{(v\Delta t)^2}{d_1^2} - 2 \frac{v\Delta t}{d_1} \cos \theta};
\eeas 
as such, the the timing error is a function
of the time interval and the velocity, as well as
the local elevation for the rover.

These examples demonstrate why in scenarios 
when both the clock parameters
and the position of the rover need to be estimated, it is
advantages to use available observables to estimate {\em both}
if we are using ranging to create time observables.

A few notes: 
\begin{itemize}
\item When the rover clock is accurate and
the position is unknown,
then we can estimate position from range measurements,
using at least three independent directions.
However, if we do not get pseudorange observables
at the same instance, then we need to use
an inertial measurement unit to adjust for position updates.
That is, in this case we can reduce the
GPS/Kuiper SV navigation solution by one variable, and use
three pseduorange observables to estimate the
rover position, when the rover is stationary.
%
When the rover moves, then we can update the position
every time we obtain the required observables 
from GPS or Kuiper SV network, if available.
In between receiving the observables, we can
rely on the onboard IMU to propagate the
translational states of the rover. 
Combining these propagated states with observables
when available is really what a filter does.

\item Similarly, if the position of the rover is
known, without relying on the clock model, the 
time synchronization protocol can use averaging schemes
to correct for clock offset; we saw an example of such a strategy
for CT synchronization before. We have developed 
the underlying equations to show how many
averaging steps are needed to achieve a particular
synchronization error.
\end{itemize}

Before we delve into developing nonlinear filters for
the rover position and time estimation, let us review the basic
ingredients for the corresponding measurement model
as well as the underlying ``update'' equations.
Let us start with the measurement model. 
We have been using time signaling as our baseline
observable- that is, we use time stamping as a proxy
for distance, that in conjunction with speed of signal
propagation leads to a time observable as well.
%-----

Now, let us review the update dynamics for different
modules on the rover receiver. These models are important
not only for designing filters, but predicting how
the uncertainties grow over time. We saw a discussion
on this for the CT clock dynamics before- but let us 
quickly review how that model can be adopted for the rover
clock. We then look at the translational
dynamics of the rover; these models are important
to not only developing position/timing algorithms,
but also for propagating uncertainties.

For the rover state, we consider an eight dimensional
state vector at time $t$ as, 
\beas
\bm s(t)= \leftm{c} \bm r(t)\\ \bm v(t)\\ \bm x(t) \rightm 
\eeas 
where $\bm r(t)$ is the position vector in some three-dimensional
coordinate frame (ENU, ECI, or ECEF), $\bm v(t)$ is the velocity
vector in the same coordinate, with respect to $t$,
and $\bm x$ is the clock state; for the two-state differential
model of the clock and since we will be combining position
and timing in one state vector, we let
\beas
\bm x^{1}(t): &\quad& \mbox{signal speed $\times$ clock offset at time $t$}\\
\bm x^{2}(t): &\quad& \mbox{signal speed $\times$ clock skew at time $t$}
\eeas 
in scenarios where we have multiple clocks we can
use $\bm x_i$ or even something like $\bm x_{{\mbox \tiny R}_i}$ when this augmented notation is justified.
Similarly, when we have multiple rovers, aerial vehicles, SVs,
etc., we will use the notation $\bm s_i$ or $\bm s_{{\mbox \tiny R}_i}$ to refer to the state of the node; a notation like
$\bm r^{j}_{i}$ refers to the $j$th component of the position vector
for node $i$; for position and velocity vector we can also
adopt a notation like $\bm r^{x}_{i}$, that refers to
the position of the $i$th node along the $x$-axis of the implied coordinate frame. For the clock states we also use
the notation
\beas
\bm x^{\phi}(t): &\quad& \mbox{signal speed $\times$ clock offset at time $t$}\\
\bm x^{f}(t): &\quad& \mbox{signal speed $\times$ clock skew at time $t$}
\eeas 
to denote the clock phase offset and skew in distance units;
the two-state model is depicted in Figure~\ref{2-state-clock}.
\begin{figure}
\begin{center}
\scalebox{0.04}{\includegraphics{eps/clockeps/2-state-clock-model.jpg}} 
\caption{2-state clock model}\label{2-state-clock}
\end{center}
\end{figure}
Lastly, note that the state of the node can also be augmented to capture {\em other} states of the nodes, for example the orientation of the vehicle or its antenna.

Now, let us consider the evolution of the state in 
discrete time, that is, we consider a sampling
interval $\Delta t$, and examine (instead of $\bm s(t)$),
the evolution of the state at sampled time steps,
\beas
\bm s(k \, \Delta t)  \leadsto \bm s((k+1) \Delta t) ;
\eeas 
in a nutshell, this evolution is dictated by memory,
controlled input, and exogenous inputs. 
Let us consider first the clock
states- we have discussed this to some extent
in the previous chapters but provide an overview
(and more inline with its subsequent integration with 
translational states) here.

Let us first note that the continuous time clock model
assumes the form
\beas
\leftm{c} \dot{\bm x}^{\phi}(t)\\ \dot{\bm x}^{f}(t) \rightm = \leftm{cc} 0 & 1 \\ 0 & 0 \rightm \leftm{c} {\bm x}^{\phi}(t)\\ {\bm x}^{f}(t) \rightm + \leftm{c} \bm w_c^{\phi}(t) \\ \bm w_c^{f}(t) \rightm;
\eeas
where $\bm w_c^{\phi}(t)$ and $\bm w_c^{f}(t)$ are
zero-mean white processes that capture the phase and
frequency noise statistics, assumed to be uncorrelated.
The covariance of these noise terms
\beas
Q_c(t) = \mathbb E \{ \bm w_c(t) \bm w_c(t)^{\top} \} = \leftm{cc} q_c^{\phi & 0 \\ 0 & q_c^{f}} \rightm
\eeas 
\beas
Q_d(k) = \mathbb E \{ \bm w_d(k) \bm w_d(k)^{\top} \}
\eeas
where
\beas
Q_d(k) = \leftm{cc} q_c^{\phi} \Delta t + q_c^{f} \Delta t^3/3 & q_c^{f}\Delta t^2/2 \\ q_c^{f} \Delta t^2/2 & q_c^{f} \Delta t \rightm
\eeas
for the resulting discrete time dynamics assumes the form,
\beas
\leftm{c} {\bm x}^{\phi}(k+1)\\ {\bm x}^{f}(k+1) \rightm = \leftm{cc} 1 &\Delta t \\ 0 & 1 \rightm \leftm{c} {\bm x}^{\phi}(k)\\ {\bm x}^{f}(k) \rightm + \leftm{c} \bm w_d^{\phi}(k) \\ \bm w_d^{f}(k) \rightm;
\eeas

Now, let us see how this model would fit the typical
way of specifying frequency stability for a clock.
A popular specification is in terms of the
Allan variance, capturing the RMS value
in the variations of the the clock frequency over 
{\em a particular time interval}. As Figure~\ref{allan-variance}
depicts, these specification is really time dependent.
\begin{figure}
\begin{center}
\scalebox{.9}{\includegraphics{eps/clockeps/allan-variance.jpg}} 
\caption{Square root of Allan variance (or Allan deviation) plot with asymptotes for a typical ovenized crystal oscillator- include the reference}\label{allan-variance}
\end{center}
\end{figure}

So let us see how such a specification maps to the
covariance in the discrete model of the clock.
As Figure~\ref{allan-variance} shows, the slope
of the variations are typically in terms of the
Allan variance parameter, say, 
$h_{0}$, $h_{-1}$, $h_{-2}$.
% \newline
% {\color{red} \bf Question:} {the typical sampling time for viewing a Kuiper SV from a mobile unit?}
% \newline

The clock state model uses the clock parameters
specified for the clock, in particular,
what are know as Allan variance parameters, 
$h_{0}$, $h_{-1}$, $h_{-2}$, etc; these are
depicted in Figure~\ref{allan-variance}.
The table below shows some typical values
for various types of clocks.
\begin{center}
\begin{tabular}{c|c|c|c}
{\bf Timing standard} & $h_0$ & $h_{-1}$ & $h_{-2}$
\\
\hline \hline
TCXO (low quality) & 2 \times 10^{-19} & 7 \times 10^{-21} & 2 \times 10^{-20}\\
TCXO (high quality) & 2 \times 10^{-21} & 1 \times  10^{-22} & 3 \times 10^{-24}\\
OCXO & 2 \times 10^{-25} & 7 \times  10^{-25} & 6 \times 10^{-25}\\
Rubidium & 2 \times 10^{-22} & 4.5 x 10^{-26} & 1 \times 10^{-30}\\
Cesium& 2 \times 10^{-22} & 5 \times  10^{-27} & 1.5 \times 10^{-33}\\
\caption{Typical Allan Variance}
\end{tabular}
\end{center}
Now, we look at the entries of the covariance matrix $Q_d$,
we can make the correspondence,
\beas
\mathbb E \{ (\bm x^\phi)^2\} &=& q_c^{\phi} \Delta t + q_c^{f} \Delta t^3/3\\
\mathbb E  \{ \bm x^f \, \bm x^\phi \}&=& q_c^{f} \Delta t^2/2 \\
\mathbb E  \{ (\bm x^f)^2 \} &=& q_c^{f} \Delta t,
\eeas 
over sample interval $\Delta t$; the firs expression also implies
that the average RMS value for $\bm x^\phi$ is
\beas
\sqrt{\frac{q_c^{\phi}}{\Delta t} + \frac{q_c^{f} \Delta t}{3}}
\eeas
There are a few procedures to go from clock specification
to the covariance of the dynamics. The first one is
using the {\em frequency stability} often specified in terms
of parts per million (ppm). This is the RMS value of
the frequency variation, so this would correspond to having
\beas
(c \times \mbox{frequency stability})^2 &\approx& \mathbb E  \{ (\bm x^f)^2 \}/\Delta t^2 \\
&=& (Q_d)_{22}/\Delta t^2
\eeas 
hence, for a frequency stability value of $10^{-7}$, we have
\beas
q_c^{f}/\Delta t = (c \times \mbox{frequency stability})^2
\eeas 
and for a system like GPS with a sampling interval
$\Delta t = 1$ sec, we end up with 
\beas
q_c^{f} = (c \times \mbox{frequency stability})^2
\eeas
Now one can argue that the the phase variation term,
the contribution the two terms should be equal and
as such, 
\beas
q_c^{\phi} \Delta t = q_c^{f} \Delta t^{3}/3
\eeas
and therefore,
\beas
q_c^{\phi} = q_c^{f} \Delta t^2/3
\eeas 
so again for $\Delta t =1$ sec, we end up with 
\beas
q_c^{\phi} = q_c^{f}/3
\eeas 
Hence, if we are given frequency stability of $10^{-7}$ say
over 1 second, then we can extrapolate to construct
the covariance matrix for the desired duration- hence,
for a $\Detal t=100$ sec (what I have used in the simulations),
then,

\beas
Q_d(k)= \leftm{cc} 3.33\times 10^{-09} & 5.00 \times 10^{-11}\\
 5.00\times 10^{-11} & 1.00\times 10^{-12} \rightm
\eeas 
Let us see how the clock simulation would work for
the corresponding dynamic model- the frequency variation
as well as the variation of the bias over one period
is shown in,
\begin{figure}
\begin{center}
\scalebox{.5}{\includegraphics{eps/mobilityeps/data_freq_variations.png}} 
\caption{Clock frequency variation from the dynamic clock model for frequency stability of $10^{-7}$}\label{data_freq_variations}
\end{center}
\end{figure}
Now suppose that we look at the RMS value of this variation-
in principle, this RMS value should be consistent with the frequency stability value.
The RMS values as a function of time horizon the looks like,
\begin{figure}
\begin{center}
\scalebox{.5}{\includegraphics{eps/mobilityeps/freq_variations_horizon.jpeg}}
\caption{Frequency variation as a function of the time horizon for frequency stability of $10^{-7}$}\label{horizon-freq-variations}
\end{center}
\end{figure}
indicating that the frequency stability specification is valid
over a time horizon that is consistent with the simulation interval.
Another way to build the covariance is directly for
the Allan variance parameter; for example, it turns
out that 
\beas
q_c^{\phi} \approx h_0/2 \quad \mbox{and} \quad q_c^{f}\approx 2 \pi^2 h_{-2}
\eeas
However, the important point is that in order
to build the covariance from Allan variance, we need
to specify the time interval for which the variance 
is being measured.

Now let us look at the clock bias- the essence of the
the so-called extended (Kalman) filtering
scheme is to estimate this bias and simultaneously estimate
the position.


\begin{figure}
\begin{center}
\scalebox{.5}{\includegraphics{eps/mobilityeps/data_bias_variations.png}}
\caption{Rover clock bias variation as a function of the time horizon for frequency stability of $10^{-7}$}\label{horizon-freq-variations}
\end{center}
\end{figure}





Now suppose that we consider the time and position accuracy
as a function of the clock frequency stability- the



Using this, we can then obtain a complete picture
for parameterizing the evolution,
\beas
\bm x(k+1) = \leftm{c} \bm x^\phi(k+1)\\\bm x^f(k+1) \rightm=\leftm{cc} 1 & \Delta t\\
          0 & 1 \rightm \leftm{c} \bm x^\phi(k)\\\bm x^f(k) \rightm  + \leftm{c} \bm w^\phi(k)\\\bm w^f(k)\rightm;
\eeas
essentially, the Allan variance for the clock, dictates
the covariance of the noise term in the evolution equation.
Hence, one of the main issues in the Kuiper network,
consisting of various nodes (CT, SV, mobile, aerial)
is that are developing a synchronization protocol
to achieve time synchronization for clocks with
distinct noise covariances.

Similarly, we can develop a dynamic model
of the position and velocity of the vehicle,
as well as the combined model with and without the IMU.
The corresponding block diagrams are shown
in Figures~\ref{2-state-position} and~~\ref{2-state-position-imu}.

\begin{figure}
\begin{center}
\scalebox{0.04}{\includegraphics{eps/clockeps/2-state-position-model.jpg}} 
\caption{2-state rover dynamics model}\label{2-state-position}
\end{center}
\end{figure}

\begin{figure}
\begin{center}
\scalebox{0.04}{\includegraphics{eps/clockeps/2-state-position-model-imu.jpg}} 
\caption{2-state position model with IMU and GPS/Kuiper integration}\label{2-state-position-imu}
\end{center}
\end{figure}


Now, let us go back the state equation 
\beas
\bm s(k+1) = A \bm s(k) + \bm w(k)
\eeas 
where the state was defined previously and
the $A$ matrix is of the form
\beas
\leftm{cc|cc} I & \Delta t \, I & 0 & 0\\
          0 & I       & 0 & 0\\
          \hline
          0 & 0 &    1 & \Delta t\\
          0 & 0 &    0 & 1 \rightm
\eeas
where $I$ denotes the $3 \times 3$ identity matrix;
hence the upper left block is a $6 \times 6$ matrix.
This equation can be used to propagate the
position and velocity of the vehicle
as well as its clock. 

For the rover state, we consider an eight dimensional
state vector at time $t$ as, 
\beas
\bm s(t)= \leftm{c} \bm r(t)\\ \bm v(t)\\ \bm x(t) \rightm 
\eeas 
where $\bm r(t)$ is the position vector in some three-dimensional
coordinate frame (ENU, ECI, or ECEF), $\bm v(t)$ is the velocity
vector in the same coordinate, with respect to $t$,
and $\bm x$ is the clock state; for the two-state differential
model of the clock and since we will be combining position
and timing in one state vector, we let
\beas
\bm x^{1}(t): &\quad& \mbox{signal speed $\times$ clock offset at time $t$}\\
\bm x^{2}(t): &\quad& \mbox{signal speed $\times$ clock skew at time $t$}
\eeas 
in scenarios where we have multiple clocks we can
use $\bm x_i$ or even something like $\bm x_{{\mbox \tiny R}_i}$ when this augmented notation is justified.
Similarly, when we have multiple rovers, aerial vehicles, SVs,
etc., we will use the notation $\bm s_i$ or $\bm s_{{\mbox \tiny R}_i}$ to refer to the state of the node; a notation like
$\bm r^{j}_{i}$ refers to the $j$th component of the position vector
for node $i$; for position and velocity vector we can also
adopt a notation like $\bm r^{x}_{i}$, that refers to
the position of the $i$th node along the $x$-axis of the implied coordinate frame. For the clock states we also use
the notation
\beas
\bm x^{\phi}(t): &\quad& \mbox{signal speed $\times$ clock offset at time $t$}\\
\bm x^{f}(t): &\quad& \mbox{signal speed $\times$ clock skew at time $t$}
\eeas 
to denote the clock phase offset and skew in distance units;

hence, the clock bias and skew are now expressed in terms of 
meters and meters per second.

\subsection{Extended Kalman filter: SV-mobile network}
We first consider the case of no IMUs and show how
standalone Kuiper time-stamped signals can be used
for an extended Kalman filter for estimating the clock parameters
as well as the position of the mobile unit.
%----
This is depicted in Figure~\ref{only-time-stamping}.
%
\begin{figure}[ht]
\begin{center}
\scalebox{0.05}{\includegraphics{eps/clockeps/Kuiper-mobile-3}}
\caption{Kuiper mobile unit position/timing using SV time stamping}
\label{only-time-stamping}
\end{center}
\end{figure}
As we will be considering mobile units in conjunction
with Kuiper SV dynamics, we need to setup the appropriate
coordinate frames. For example, it is not necessary convenient
to express the motion of the mobile unit in an earth-centered 
inertial (ECI) frame. For the development of the 
filtering scheme discussed below, we need three
distinct coordinate frames,
\begin{enumerate}
\item east-north-up (ENU) frame that is a Euclidean plane tangent to a point on earth's surface/
\item earth-centered inertial frame (ECI) 
\item earth-centered earth-fixed (ECEF) frame (that rotates with earth)
\end{enumerate}

The SV dynamics is usual expressed in terms of ECI frame;
the usual way that Newton's second law is combined universal gravitation
to derive the satellite-earth relative dynamics (in an inertial frame).
One can then translate coordinates in ECI frame to the ECEF. 
%
Similarly, we can write the mobile unit dynamics in ENU frame,
then translate it to ECEF frame so that we can do
the appropriate computations consistently between
the positions of the SV and the mobile unit.
These three coordinate frames are depicted in Figure~\ref{ENU-ECI-ECEF}.
%
\begin{figure}[ht]
\begin{center}
\scalebox{0.05}{\includegraphics{eps/clockeps/ENU-ECI-ECEF}}
\caption{ENU vs. ECEF vs. ECI}
\label{ENU-ECI-ECEF}
\end{center}
\end{figure}

Now let us consider the ground mobile unit traversing
a particular trajectory; the coordinates of the mobile
unit are expressed conveniently in the ENU coordinates.
\begin{figure}[ht]
\begin{center}
\scalebox{0.05}{\includegraphics{eps/clockeps/4-sv-mobile-measurements}}
\caption{Measurement geometry in a Kuiper time-stamping scheme has implications about the position/timing estimation errors}
\label{4-measurements}
\end{center}
\end{figure}

For our working examples, we have considered a mobile unit
that traverses a few types of trajectories- for example,
to the east, then up, then west, and down the original location with a given
speed. For now, we have let the time interval for completing
this trajectory be equal to the orbital period for
the Kuiper (phase I) SVs. We assume that we have
time-stamped measurements from certain Kuiper SVs
available at the mobile units. As pointed out
above, we do not assume having access to IMUs for now
(but they will be included in the architecture/algorithms).\\

The steps that we have followed to develop the code:
include:
\begin{enumerate}
	\item Choose the number of SVs in view for which you could receive time stamp information.
	\item We can decide which SV-mobile unit link should be used; we can use all of them or only one. The current setup uses the distance to the SV as a criteria for using the time stamp measurements from a given satellite- this is justified by the fact that delays and uncertainties induced by the troposphere are generally a function of the path length traversed by the signal. As such, choosing the closest SV can result in smaller uncertainty in estimating the path delays. On the other hand, as we discussed before, choosing directions from mobile unit to SV that are ``independent'' result in better conditioning and estimation performance for the position and timing for the mobile unit.
	\item The time stamped signal allows us to construct the pseudorange, essentially the ranges between the mobile units and the SVs, that also contains information on the mobile unit {\em clock offset}.
	\item The mobile unit receiver now uses these constructed pseudorange observables in the context of an extended Kalman filter (EKF).
	\item The EKF input hence is a sets of noisy measurements of the position of the mobile unit with respect to the SV. EKF has to make do a few coordinate transformations to make sure that the variables are in the right coordinate frame. The reference to ``extended'' in the EKF is due to the fact that the range measurements are nonlinear function of the mobile unit position- as such, we have to use a recursive linearization scheme in order to combine pseudorange measurements with the knowledge of the translational and clock dynamics.
	\item The EKF then updates the temporal and spatial coordinates of the mobile units via estimating the clock offset and skew. This is done by having an internal model of the mobile unit translational (position and velocities in the ENU coordinate frame) as well as the clock dynamics- there is flexibility in making the translational and clock dynamic more elaborate- including updating covariance representation of the clock bias/skew noise due to environmental effects like temperature.
	\item The covariances for position and timing errors are generally a function of: 1) number of SV time stamp measurements, 2) directions of measurements for the pseudo range measurements, 3) the error in estimating the path delays to the SV due to troposphere, 4) quality of the clock onboard the mobile unit, 5) the balance between speed of the mobile unit and the interval between two time stamp updates from the SVs- this is a more involved analysis- essentially, when we consider EKF (or more generally, any recursive filtering) we often assume regularly sampled updates- when we consider intermittent sampling, then even convergence/boundedness of the recursive algorithms is tricky.
\end{enumerate}

Using the above setup, we then proceeded to examine how various uncertainties
and clock parameters and delay uncertainties in the signal 
propagation, as well as mobility, effect the synchronization
performance on the mobile unit. Below we examine some of the
simulation results in regards to assessing the effect of,
\begin{enumerate}
	\item clock quality, temperature variations
	\item number of Kuiper SVs in view
	\item path delay uncertainty due troposphere 
	\item update rate 
\end{enumerate}
on the position and time estimation for the Kuiper 
mobile units using time stamping signal from the SVs.
\begin{shaded*}
The simulation results/plots presented in this chapter for
ground mobile units have been generated by the 
python script "mobile-unit-sync.ipynb"
\end{shaded*}¬j

\begin{figure}
\vspace{-.1in}
\begin{center}
\includegraphics[width=0.4\textwidth]{eps/clockeps/rover-clock_4sv.png}
\includegraphics[width=0.4\textwidth]{eps/clockeps/rover-estimate1.png}
\end{center}
\caption{Using time-stamped signals by accessing 4 SVs for time synchronization and positioning} \label{pnt-rover-time-stamping}
\end{figure}


\begin{figure} %[8]{r}{0.3\textwidth}
\begin{center}
\includegraphics[width=0.4\textwidth]{eps/clockeps/rover-estimate_2svs.png} \quad
\includegraphics[width=0.4\textwidth]{eps/clockeps/rover-estimate_2sv_acc.png}
\caption{Using 2 SV time-stamped signals for rover positioning and time synchronization without an IMU (left panel) and using a simple model of an accelerometer (right panel)} \label{pnt-2-svs}
\end{center}
\end{figure}
{One of the observations so far has been that using EKF,
the time synchronization can be achieved efficiently by relying
only a few (even one) SV time stamped signals, but the position estimates become
progressively worse by relying on less satellites.}

This last observation motivated looking at how we can fuse information
from the inertial measurement units, like an accelerometer, can
be used to future improve the positioning of the mobile unit using
fewer satellites. The IMU modeling takes a bit more effort,
however, the initial assessment (using a simple model of including
a noisy velocity information in the EKF) looks promising.
If desired, we can proceed to,
\begin{itemize}
	\item develop a more IMU dynamic model for Kuiper model unit and assess the corresponding  error statistics,
	\item develop an EKF that integrates time stamping from Kuiper SVs and the IMUs using more analysis on the corresponding error statistics.
\end{itemize}
Another possibility to examine if we can use doppler shift as a proxy
for velocity information that can then been integrated in the estimation/filtering
process.
%
% \subsection{Geolocation}

% using a uav to see localize 
% the rover.

\subsection{Using IMUs on the Rover for Positioning and Timing}
%---
As we discussed in the previous section, one 
can use time-stamped signals from Kuiper SVs
at the rover's receiver to resolve for  positioning
and time synchronization. However, in situations
where the SVs are not accessible, the rover
can rely on its onboard sensors to
propagate the position and velocity of the rover,
until the next time-stamped signal from Kuiper SV (or GPS when available) becomes available.\footnote{Note that we can also develop this setup for indoor rovers using ``time'' anchors in the building- another situation would be in urban areas with tall buildings that can periodically block 
GPS availability.}
%----
A prime example of such an onboard
sensor is an {\em inertial navigation unit} (IMU), consisting
of accelerometers and gyroscopes. We will discuss IMUs
shortly, as they become indispensable for
positioning and time synchronization for
aerial vehicles as well. Other types of data that can
be used for positioning of the rover, include
odometry (using wheel sensors to measure the wheel rotation),
electronic compass (for heading),
cameras, sonar (for indoor application),
and LIDAR, and digital maps.\footnote{A potential direction for further work is to explore augmenting the proposed position/time synchronization for Kuiper units with these other sensing/measurement modalities.}

The overall block diagram for how we envision 
integrating IMUs for the position and time synchronization is shown in Figure~\ref{ins-integration};
by inertial navigation system (INS), one
often means the combination of inertial measurement
units as well as the necessary
post-processing (known
as {\em mechanization} in GPS literature) 
that leads to an estimate of 
position and velocity of the rover.
\begin{figure}[h]
\begin{center}
\includegraphics[width=0.7\textwidth]{eps/clockeps/integrated-ins-kuiper}
\end{center}
\caption{Integration of inertial navigation system (INS) with Kuiper SV time-signaling (and potentially GPS)} \label{ins-integration}
\end{figure}

Let us start with the basics setup
of an accelerometer
for estimating the position of the rover based on an onboard sensors. As its name suggests,
an accelerometer estimates the acceleration experienced
by the rover- in its basic form, 
this is done by ``measuring'' 
the displacement of a proof mass inside the
accelerometer casing.\footnote{For MEMS accelerometers, the idea is similar yet the realization of the ``proof mass'' is more creative.}
%---
In reference to Figure~\ref{accel}, suppose that
the position of the center of mass of the proof 
mass is denoted by $x(t)$ and that of the casing
as $y(t)$. We assume that the proof mass has mass $m$ 
and is attached to the body of the casing by springs
with spring constants $k/2$.
\begin{figure}[h]
\begin{center}
\includegraphics[width=0.2\textwidth]{eps/clockeps/accelerometer}
\end{center}
\caption{Simplified model of an accelerometer} \label{accel}
\end{figure}
We can now represent the dynamics of the proof mass 
by the second order differential equation of the form,
\beas
m \, \frac{d^2}{dt^2} x(t)= m \, \ddot{x}(t) = k \, \delta(t) ,
\eeas 
where $\delta(t)= y(t)-x(t)$ is the relative displacement
of the casing with respect to the proof mass and
the time derivative is with respect to an inertial frame;
as such,
\beas
\ddot{x}(t) = \frac{k}{m} \delta(t).
\eeas 
Hence the displacement of the proof mass
becomes a proxy for the acceleration experienced
by the proof mass, and by extension, the
acceleration experienced by the casing. 
In fact, one can compute the transfer
function between the displacement of
the proof mass and the chasing as
\beas
\frac{y(j\omega)}{x(j\omega)} = \frac{1}{1 - (m/k) \omega^2},
\eeas 
so when $\omega \ll \sqrt{k/m}$,
the displacement of the proof mass is
accurate indicator of the acceleration
of the body. The displacement of the
proof mass, in the meantime, can be 
picked up by a capacitive transducer, 
converting the displacement
to a voltage (MEMs accelerometer).

In a nutshell, the accelerometer's output of
acceleration $a_{\mbox{\tiny meas}}(t)$ in one
axis is a generally a biased/noisy version of the
proof mass acceleration $a(t)$;
as such we can write, 
\beas
a_{\mbox{\tiny meas}}(t)= K_{\mbox{\tiny acc}} a(t) + b_{\mbox{\tiny acc}} + \zeta_{\mbox{\tiny acc}}
\eeas
where $K_{\mbox{\tiny meas}}$ is the gain of
the accelerometer, $b_{\mbox{\tiny acc}}$ is its
bias, and $\zeta_{\mbox{\tiny acc}}$ is the 
accelerometer noise.
%%%%
This would provide a noisy-biased acceleration measurement in one axis; note that the way we have setup the problem so far, all directions are aligned- the casing and the proof mass movements are aligned with the inertial $x$-axis.

Before we delve into a more generation situation,
consider the accelerometer assembly in a configuration
shown in Figure~\ref{general-config-accel}.
\begin{figure}
\begin{center}
\includegraphics[width=0.5\textwidth]{eps/clockeps/strapdown}
\end{center}
\caption{Assembling accelerometers and gyroscopes in one unit} \label{general-config-accel}
\end{figure}
Again, assuming that all the corresponding coordinate frames are aligned, 
the output of the accelerometer assembly can
be represented in vector form as,
\beas
\bm a_{\mbox{\tiny meas}}(t) = A \bm a(t) + \bm b_{\mbox{\tiny acc}} + \bm \zeta_{\mbox{\tiny acc}} ,
\eeas 
where the variables are now the ``vector'' of the three
acceleration terms in each inertial axis; $A$ in this
case is a $3 \times 3$ matrix. 
Once acceleration of the casing has been
estimated, we can then proceed to estimate the position 
of the rover with respect to its initial position
and velocity by integration, as
\beas
\bm v_{\mbox{\tiny est}} (t) = \bm v (t_0) + \int_{t_0}^{t} \bm a_{\mbox{\tiny meas}}(t) \, dt \quad \mbox{and} \quad \bm r (t) = \bm r (t_0) + \int_{t_0}^{t} \bm v_{\mbox{\tiny est}} (t) \, \dt.
\eeas 
Note that the acceleration measurements as well as
the integration would be approximate/noisy, and hence,
these estimates are combined with those obtained
from direct Kuiper SV time-stamping; for example,
an architecture such as that shown in Figure~\ref{loose-1}
can be used to fuse the position/velocity
estimates from each channel.
\begin{figure}[h]
\begin{center}
\includegraphics[width=0.7\textwidth]{eps/clockeps/loosely-coupled} 
\end{center}
\caption{Integration of Kuiper SV time-stamping and IMU; between two updates from the Kuiper SV, the IMU integrates acceleration forward to estimate the position and the velocity of the
vehicle before the next Kuiper time-stamped signal has been received. The Kalman filter then combines the IMU estimation of position/velocity, with the measurements received from the satellite, to obtain a statistically optimal estimate of the rover's position, velocity, and timing error.} \label{loose-1}
\end{figure}

In Figure~\ref{loose-1}, we have made a reference to ``Mechanization'' of the IMU measurement- hence let us explain what this means. Mechanization refers to the process of converting IMU outputs to states that
are of interest to us in the {\em right coordinate frame.} For example, consider again the ENU frame- or the {\em local navigation} frame $n$
(Figure~\ref{ENU-ECI-ECEF-repeat}), where we are interested to estimate the position and velocity of the vehicle.
\begin{figure}[ht]
\begin{center}
\scalebox{0.05}{\includegraphics{eps/clockeps/ENU-ECI-ECEF-2}}
\caption{ENU vs. ECEF vs. ECI}
\label{ENU-ECI-ECEF-repeat}
\end{center}
\end{figure}
What this means that we need the acceleration 
of the accelerometer casing expressed in 
this frame (such frames are often called ``navigation frames''). Hence if we denote the position of accelerometer casing/rover in the (navigation)
ENU frame as
\beas
{\bm r}_{\mbox{\tiny R}/n}
\eeas 
then in order estimate this position via
the integration of the accelerometer output, we need
the output of the accelerometer to be converted
to
\beas
{\bm a}_{\mbox{\tiny R}/n};
\eeas 
when it is implicit that we are referring to the rover we simply
the above notation, e.g., ${\bm r}_{n}$ is the position of the rover
in the frame of the navigation frame $n$.
%---
This is where mechanization comes in as
what is measured by the accelerometer is {\em not}
${\bm a}_{\mbox{\tiny R}/n}$!
%
In order to understand what is going on, consider
again the accelerometer casing now attached to the rover whose movement is monitored in the ENU frame ${\cal F}_n$. 
% In order to streamline the notation
% in this section, when we adopt the convention
% \beas
% {\bm p}_{\mbox{\tiny R}} = {^{\ell}{\bm p}}_{\mbox{\tiny R}}
% \eeas
% for any vector $\bm p$.
We assume that the three accelerometers in the
IMU are aligned with the rover's axes and located at the center of mass of the vehicle. We call this
the rover's frame ${\cal F}_{\mbox{\tiny R}}$; note that
if we wanted to explicitly include the antenna frame, we have
to include a conversion as the antenna will typically not be
at the vehicle center of mass.

Before we derive the relations behind the
mechanization process, let us include a quick note on 
coordinate transformations.
Suppose we have two coordinate frames ${\cal F}_{1}$ and
${\cal F}_{2}$ and they are related to each
other via a coordinate transformation $\cal C$ (or rotation):
\beas
{\cal C}_{1}^{2} \quad \mbox{transformation required on ${\cal F}_{1}$ to obtain ${\cal F}_{2}$;} \quad \mbox{or, orientation of ${\cal F}_{2}$ with respect to ${\cal F}_{1}$}
\eeas 
This also means that if we have a vector $\bm p$
that has a representation $\bm p^{{\cal F}_{1}}$, in order to obtain its representation in ${\cal F}_{2}$, we can do,
\beas
\bm p^{{\cal F}_{2}}= {\cal C}_{1}^{2} \; \bm p^{{\cal F}_{1}}
\eeas 
The convention is that this transformation (rotation) expresses the rotation
of the first frame to the second in the {\em right-handed} coordinate. That is for the 2D case, the rotation of the frame is counter-clockwise. 

Having this notation in place, let us revisit 
understanding what quantity the accelerometer (whose axes are aligned with the rover coordinate frame) actually measures. In a nutshell, the accelerometer measures
the ``specific force" in the body frame, where
the specific force is the net inertial acceleration.
Suppose 
\beas
{\bm r}_{\mbox{\tiny R}/n}
\eeas 
is the position of the rover in the ENU frame ${\cal F}_{n}$.
Then its inertial velocity and acceleration would
be
\beas
\frac{d}{dt_i} \, {\bm r}_{\mbox{\tiny R}/n}(t) \quad \mbox{and} \quad \frac{d^2 }{dt_i^2} \, {\bm r}_{\mbox{\tiny R}/n}(t),
\eeas 
where the significant of ``$i$'' as a subscript in the denominator is that the time derivative is
taken with respect to the inertial frame.

The acceleration that the accelerometer 
measures then is 
\beas
\{\frac{d^2}{dt_i^2}\, {\bm r}_{\mbox{\tiny R}}(t)\}^{\mbox{\tiny R}}
\eeas 
where ``${\mbox{\tiny R}}$'' is the abbreviated for ${\cal F}_{{\mbox{\tiny R}}}$ (the rover--or the casing--body frame).
%
Now mechanization gets to work! That is, we need
to convert 
\beas
\{\frac{d^2}{dt_i^2} \, {\bm r}_{\mbox{\tiny R}}(t)\}^{\mbox{\tiny R}} \leadsto \{ \frac{d^2 }{dt_{n}^2}  \, \; {\bm r}_{\mbox{\tiny R}}(t)\}^{n} 
\eeas 
since it is the double integration of the quantity on
the right hand side that leads to
\beas
{\bm r}_{\mbox{\tiny R}}^n (t)
\eeas 
that should be read as the position of the rover 
in the ENU coordinate frame resolved in the 
ENU coordinate frame!

The key relations that allows us to do this
conversion are as follows: consider 
a vector $\bm p$ and its derivatives with respect to
coordinate frames ${\cal F}_1$ and ${\cal F}_2$;
then 
\beas
\frac{d}{dt_1} \bm p= \frac{d}{dt_2} \bm p + \bm \omega_{2/1} \times \bm p 
\eeas
where $\bm \omega_{2/1}$ is the angular acceleration
vector between the two frames, i.e., rate at which
${\cal F}_2$ rotates with respect
to ${\cal F}_1$.
%---
Note that in the above
relation we have not included where these vectors are
being resolved in.
When we continue this for the second derivative, we obtain
\beas
\frac{d^2}{dt_1} \bm p= \frac{d^2}{dt_2} \bm p+ 2 \, \bm \omega_{2/1} \times \frac{d}{dt_2} \bm p + \bm \omega_{2/1} \times (\bm \omega_{2/1} \times \bm p).
\eeas
This is a key relationship for the mechanization process, as the acceleretion measured by the accelerometer is the {\em total acceleration} that involves not only gravity but also acceleration induced by the various
rotating coordinate frames with respect to the inertial frame.


Now let us come back to the accelerometer measurements;
even in the simpler case of not dealing with the antenna frame,
we have three different coordinate frames: the
ECI frame that is inertial (this will be denoted by ${\cal F}_i$), then ECEF frame that is rotating with earth (denoted by ${\cal F}_{\mbox{e}}$, then the ENU frame (denoted by ${\cal F}_n$, and then the rover/accelerometer frame
(denoted by ${\cal F}_{\mbox{\tiny R}}$)!
%----
Adopting the above relationships, we can then
write,
\bea
\{ \frac{d^2}{dt_i^2} \; {\bm r}_{\mbox{\tiny R}}(t)} \}^{n}  = \{ \frac{d^2}{dt_{n}^2} \; {\bm r}_{\mbox{\tiny R}}(t)\}^{n} + 2 \, \{ \bm \omega_{n/i} (t)\}^n \times \{ \frac{d}{dt_n} {\bm r}_{\mbox{\tiny R}}(t) \}^n
+ \{ \bm \omega_{n/i}(t) \times (\bm \omega_{n/i} (t) \times {\bm r}_{\mbox{\tiny R}}(t)) \}^{n} \label{total-accel}
\eea

In (\ref{total-accel}), the second term on the right hand side of the equation is called Coriolis acceleration and the last terms of the right hand side is called centrifugal acceleration.
The difference between these two accelerations,
one that is essentially measured by the accelerometer,
and one that is needed to estimate the position of the
rover in the local frame, it is important to know where the frame is. For example, consider the magnitude of
\beas
\| \{ \bm \omega_{n/i}(t) \times \bm \omega_{n/i} (t) \times {\bm r}_{\mbox{\tiny R}}(t) \}^{n} \|
\eeas 
and assume that $\bm \omega_{\mbox{\tiny e}/i} \approx \bm \omega_{\ell/i}$ which is about $7 \times 10^{-5}$ rad/s.
Hence the centerifugal force becomes a function of
the position of the rover in the local frame.
The Coriolos acceleration on the other hand 
is a function of the speed of the vehicles.


Now, we use the setup shown the figure below, as our blueprint,
to show precisely how the the measurement from
the accelerometer can be used to estimate the position
of the rover.
\begin{figure}[h]
\begin{center}
\includegraphics[width=0.66\textwidth]{eps/clockeps/mechanization.jpg}
\end{center}
\caption{Mechanization process for an IMU} \label{mechanization}
\end{figure}

\begin{comment}
The first under ``accelemeter output'' is in fact,
\beas
{\cal R}^{\mbox{\tiny R}}_{\ell} \, ( \{ \frac{d^2}{dt_i^2} \; \bm r_{\mbox{\tiny R}} \}^\ell - \bm g^\ell )
\eeas 
where as before ${\cal F}_\ell$
and ${\cal F}_{\mbox{\tiny R}}$ are respectively
the ENU and the rover frames.

For Kuiper we can have a series of ENU (local) frames all over
the world, defining the ``local'' coordinate frames.
For example, CT terminals or GW sites can be
used to defined the origins of these coordinate frames.
%--
In Figure~\ref{mechanization}, by ``orientation'' we refer to the fact that we assume the output from the gyroscope (in the IMU) that measures
\beas
{\cal R}^{\mbox{\tiny R}}_{i}, 
\eeas 
to estimate other rotations such as
${\cal R}^{\mbox{\tiny R}}_{\ell}$ and
${\cal R}^{\mbox{\tiny R}}_{\ell}$,
since ${\cal R}^{\mbox{\tiny e}}_{i}$ and
${\cal R}_{\mbox{\tiny e}}^{\ell}$ are known.
With respect to Figure~\ref{mechanization},
the output of the accelerometer, in conjunction
with the knowledge of ${\cal R}_{\ell}^{\mbox{\tiny R}}$
can be used to find $\{\frac{d^2}{dt_i^2} \bm r\}^{\ell}$
and by extension $\{\frac{d^2}{dt_i^2} \bm r\}^{\ell}$.
This quantity is now exactly the quantity
on the left hand side of (\ref{total-accel}).
Since
\beas
\frac{d}{dt_{i}} \bm r_{\mbox{\tiny R}} = \frac{d}{dt_{\ell}} \bm r_{\mbox{\tiny R}} + \bm \omega_{\ell/i} \times \bm r_{\mbox{\tiny R}}
\eeas 
the output of the accelerometer can related to the 
acceleration in the ENU (local) frame in the form of
\beas
\bm a_{\mbox{\tiny meas}}(t) = A \bm a(t) + \bm b_{\mbox{\tiny acc}} + \bm \zeta_{\mbox{\tiny acc}} ;
\eeas 
% we now have a better appreciation for
% how the gain and the bias terms arise.
% %---
% noting that
% \beas
% \{ \frac{d^2}{dt_i^2} \bm r_{\mbox{\tiny R}} \big|_{\mbox{\tiny ENU}} = {\cal R}_{\mbox{\tiny ENU}/\mbox{\tiny ECEF}} \, {\cal R}_{\mbox{\tiny ECEF}/i} \, \frac{d^2}{dt_i^2} {^{\mbox{\tiny{ENU}}}}\bm r_{\mbox{\tiny R}} \big|_{i}
% \eeas 

% Now we have all the pieces: the flow chart from the accelerometer output looks something like this:
\end{comment}


\begin{figure}[ht]
\begin{center}
\includegraphics[width=0.7\textwidth]{eps/clockeps/loosely-coupled-1} 
\end{center}
\caption{Integration of Kuiper SV time-stamping and IMU; we pass the timing information from the top channel for the integration at the bottom channel} \label{loose}
\end{figure}

\begin{comment}
The gyroscope measures the angular velocity
of the body frame with respect to the inertial frame,
expressed in the body frame. 
\beas
\bm \omega_{\mbox{\tiny R}/\mbox{\tiny ENU}} \big|_{\mbox{\tiny R}} = {\cal R}_{\mbox{\tiny R}/\mbox{\tiny ENU}} \, (\bm \omega_{\mbox{\tiny ENU}/i} \big|_{\mbox{\tiny ENU}} + \bm \omega_{\mbox{\tiny ENU}/\mbox{\tiny R}} \big|_{\mbox{\tiny ENU}}) + \bm \omega_{\mbox{\tiny R}/\mbox{\tiny ENU}}\big|_{\mbox{\tiny R}}
\eeas 
\end{comment}

We notice that what we need is the acceleration of the
casing in some ``world" coordinate frame but
the measured acceleration is in the frame of
the casing (so-called body frame). Hence, there are
post processing steps to convert these measured accelerations to those
that can be used for integration and then fused
with those obtained using time-stamping from Kuiper SV
pseudoranges- this processing (that involves the use
of orientation and angular rate estimates
is often referred to as {\em mecahnization} in the INS/GPS integration).
Hence, if Kuiper time-stamped signaling is used for
position/velocity estimation, one has to devise a 
Kuiper mechanization process. The mechanization process
is sequence of steps that converts inertial accelerations
measured in the body frame, to the acceleration
of the vehicle with respect to local navigation frame 
(ENU) expressed in the body frame. For these conversions,
the IMU also estimates the angular rate and orientation 
of the casing as shown in Figure~\ref{loose}.
% The part that I can currently exploring
% is whether it would make sense to use
% the time synchronization obtained from the top channel
% for the integration at the bottom channel. That
% is, use an architecture that is
% depicted in Figure~\ref{loose}.

% since we have position and velocity
% dynamics/derivatives in the state update,
% we can measure $\dot {\bm v}$.
% so we can introduce one more variable,
% the include a noisy measurement- but i have to check how it works out.

% so the accelerometer would be a
% dynamics related to $\dot {\bm v}$.
% check tomorrow.
% \subsection{IMU Mechanization}
% Let us go over some of the processing
% that is needed to covert the acceleration
% experienced by the suspended mass in the
% accelerometer casing, to the acceleration
% on the vehicle- as we mentioned, after this
% conversion, one can integrate the signal
% to first get an estimate of
% the velocity, and then the position.
% As we will develop a model of the IMU in
% the code, we will expand on what needs to
% be implemented first.
% The part that we have to be careful is that
% it is the total acceleration

% \beas
% \leftm{c} a_x \\ a_y \\ a_z \rightm
% = \frac{d \bm v}{dt} + \bm \omega_{b/i} \times \bm v - R_v^{b} {\cal R}_{b/i} \leftm{c} 0 \\ 0 \\ g \rightm
% \eeas

% ${\cal R}_{b/i} \,{\cal R}_{b/i}^{-1} = I$

% where $R_v^{b}$ is the rotation matrix from
% the $v$-frame to the body frame.

% in general,

% \beas
% \frac{d}{dt_i} \bm p =  \frac{d}{dt_b} \bm p + \bm \omega_{b/i} \times \bm p
% \eeas

% so what is being measured is not necessary
% let me explain this for the 2-d case.


\begin{figure}[h]
\begin{center}
\includegraphics[width=0.5\textwidth]{eps/clockeps/body-inertial-2d} 
\end{center}
\caption{The role of vehicle's coordinate in what the accelerometer outputs} \label{loose}
\end{figure}
We solve the following differential equation to 
obtain the velocity 

\begin{comment}
\beas 
\dot{\bm v}^{n} = {\cal C}_{b}^{n} \bm f_{ib}^{b} - (2 \bm \omega^n_{e/i} + \bm \omega_{n/e}^n) \times \bm v^{n} + \bm g^n
\eeas

So we need to know what is ${\cal C}_{b}^{n}$.

\beas
{\cal C}_{e}^{n} =\leftm{ccc} - \sin \theta & \cos \theta & 0 \\
							  - \sin \phi \cos \theta & - \sin \hi \sin \theta & \cos \phi \\
							   \cos \phi \cos \theta & \cos \phi \sin \theta & \sin \theta \rightm
\eeas 
where $\theta$ is the local reference longitude and
$\phi$ is the local geocentric latitude.
where $\theta$ is the local reference longitude 
and $\phi$ is the local geocentric latitude

% we have also the coordinate frame conversion
% ENU frame to the body frame 
3-2-1 sequence \\
% $\psi$ Y- yaw\\
% $\theta$ P pitch\\
% $\phi$ R- roll

\beas
{\cal C}_{b}^{n} = \leftm{ccc} \sin \psi \cos \theta & \cos \phi  \cos \psi + \sin \phi \sin \psi \sin \theta & -\sin \phi \cos \psi + \cos \phi \sin \psi \sin \theta \\
		\cos \psi \cos \theta & -\cos \phi \sin \psi + \sin \phi \cos \psi \sin \theta & \sin \phi \sin \psi + \cos \phi \cos \psi \sin \theta \\
 \sin \theta & - \sin \phi \cos \theta & - \cos \phi \cos \theta \rightm
\eeas 


what is
\beas
\bm \omega^n_{e/i} \quad \bm \omega_{n/e}^n \quad 
\bm g ^{n} = \leftm{c} 0\\0\\ 9.8 \rightm
\eeas 

we can then do the integration to latitude and longitude
\beas
\dot{\phi} &=& \bm v^{n_x}/(R-h)\\
\dot{\lamda} &=& \bm v^{n_y}/(R-h)\\
\dot{h} &=& \bm v^{n_z}
\eeas


The IMU provides the angular rate $\bm \omega _{ib}^{b}$.
then we have to compute
\beas
\bm \omega_{b/n}^{b} = \bm \omega{b/n}^{b} - {\cal C}_{b}^{n}^{\top} (\bm \omega_{e/i}^n + \bm \omega_{n/e}^n)
\eeas


where the earth angular rate vector in the $n$-frame is
\beas
\bm \omega{e/i}^{n} = \leftm{c} \Omega \cos \phi \\ 0 \\ - \Omega \sin \phi \rightm
\eeas 
where $\Omega$ is scare earth rate of rotation and $\phi$ is
the geographic latitude.

\beas
\bm \omega_{n/e}^{n} = \leftm{c} v_{}
\eeas 

so let us put all together- imagine that
the accelerometer measures a particular specific force on the casing. Suppose
also that we have a trajectory for the
orientation of the vehicle with respect to the ENU frame- as show below.

we start with certain orientation, then
use the rate the update the orientation- this
would be ${\cal C}_{b}^{n}$. then we solve the differential equation to get $\bm v(t)$, which can then be integrated to get position estimation.


essentially, what happens is that we will have
a particular noise characteristics for the IMU;
hence the measurements would look like:

\beas
a(t) =
\eeas

we can also quantify how the uncertainty by using
the so-called Allan variance.
\end{comment}
%---
Suppose we consider the sample path for the rover as shown
in Figure~\ref{fig:sample-path}; this path is show in the local
navigation frame, the so-called ENU (east-north-up) frame.
\begin{figure}[ht]
\begin{center}
\scalebox{.6}{\includegraphics{eps/clockeps/rover-sample-path.png}}
\caption{Rover sample path for IMU-Kuiper integration} \label{fig:sample-path}
\end{center}
\end{figure}
As the rover traverses this path, the rover identifies 
a number of satellites in view and uses the corresponding
pseudorange measurements to updates it position estimate as well
as identify and correct for clock biases. 
For this example and ease of visualization, 
we have used the origin for the ENU frame
as the point with zero longitude and latitude. The
corresponding trajectories in the ECEF and ECI frames
are depicted in Figure~\ref{fig:sample-path-ecef}.

\begin{figure}[ht]
\begin{center}
\scalebox{.5}{\includegraphics{eps/clockeps/rover-sample-path-ecef.png}}
\scalebox{.5}{\includegraphics{eps/clockeps/rover-sample-path-eci.png}}
\caption{Rover sample path for IMU-Kuiper integration in ECEF (left panel) and ECI (right panel) frames over one orbital period} \label{fig:sample-path-ecef}
\end{center}
\end{figure}

The reason that we care about the ECI frame is that it is underlying
inertial frame; as such, the accelerometer (e.g., the movement of the proof mass), pertains to the acceleration of the vehicle in the ECI frame.
Following the number of steps detailed in previous section, we can
then consider the output of the accelerometer (followed by calibration and
appropriate coordinate transformations) is of the form,
\beas
\bm a_n^n(t) + \bm \zeta_{\tiny \mbox{accel}}
\eeas 
where $\bm \zeta_{\tiny \mbox{accel}}$ is the vector-valued
noise term that depends on the quality of the accelerometer; 
as before the notation the superscript in $\bm a_n^n(t)$ 
refers to the frame where the vector is expressed in,
and the subscript is which frame the acceleration has been computed
in; for example, the acceleration that is measured by the accelerometer
is really of the form,
\beas
\bm a_{\mathscr i}^\mathscr{b}(t),
\eeas 
that this, this is the inertial acceleration of the IMU housing 
expressed in the body frame; another potential notation for
this term would have been $\bm a_{b/\mathscr i}^\mathscr{b}(t)$
(acceleration of the body with respect to the inertial frame, expressed
in the body frame).

For accelerometers, and more generally IMUs, one can
quantify a measure of ``quality'' that captures the uncertainty
in the measuring acceleration of the vehicle. Now
suppose that one integrates this acceleration 
over the interval $[0,t]$ to obtain velocity,
and subsequently, position estimate.


Let us now consider how the integration of Kuiper time-signaling
and IMU would work. Let us consider the rover introduced above
and assume that the rover has access to four satellites in view
during the positioning and timing update by the rover.
The satellites' orbits in the ECI frame are show in Figure~\ref{4-sv-rover}.

\begin{figure}[ht]
\begin{center}
\scalebox{0.8}{\includegraphics{eps/clockeps/4-sv-orbits-eci.png}}
\caption{Orbit of four SV satellites in view of the rover} \label{4-sv-rover}
\end{center}
\end{figure}

In this simulation results discussed next, we used the
discretization step of $50$ seconds over the orbital period of
for Phase I Kuiper.
%--

So let us first gain some intuition on how the baseline 
positioning and timing algorithm operates. Recall
that both the clock
as well as the accelerometer drift over time- suppose
that we initialize these measurement units at time $t_0$
and start the simulation. Then (for our simulation setup),
50 seconds later, the double integration of the accelerator output
and the initial knowledge of the rover position leads to 
a new estimate of the rover position at time $t+\Delta t$,
where $\Delta t$ is the simulation/discretization step.
Now, now four SVs become available in view of the rover, 
and the rover can formalize four pseudoranges, from which
its position and timing error (clock bias) can be obtained.
The pseduorange measurements also have uncertainty associated with
them during to path delay estimates, ephemeris uncertainty, SV clock
uncertainty, etc. Now these measurements (from SV-time stamping,
accelerometers) can be combined, to obtain an estimate of the rover
position and clock bias. When we use the dynamic model of the clock
and the vehicle in order to obtain this estimate, we are essentially 
developing a integrated filter that is fusing the time-stamping
with IMU measurements. For example, in the case of 4 satellites
in view, we obtain the rover position estimation in the
ECEF frame as show in Figure~\ref{4-sv-rover-imu-estimate-fig}.

\begin{figure}[ht]
\begin{center}
\scalebox{0.8}{\includegraphics{eps/clockeps/4-sv-orbits-imu-estimate.png}}
\caption{Orbit of four SV satellites in view of the rover} \label{4-sv-rover-imu-estimate-fig}
\end{center}
\end{figure} 

A moment reflection on the role that "four" time-stamped
signals from Kuiper constellation plays reveals that these measurements are
needed to resolve for the three position coordinates and
one time coordinate associated with the clock bias-- this
is analogous to the GNSS positioning and timing.
Note that these four Kuiper SV measurements can be augmented,
or even substituted by other measurements, e.g., GPS, or
a time-position signal from a base station (or a flying aerial vehicle).
In fact, such potential augmentation schemes might allow Kuiper
to provide improved positioning and timing performance 
as compared with other LEO constellations.
In this section, the scenario that we like to consider is whether the IMU can
be used as an augmentation signal when the four SV timing
signals are unavailable.

In order to assess this, let us consider 3 SVs in view and assess
the performance of the corresponding filtering scheme for
position and timing. This is shown in Figure~\ref{3-sv-rover-imu-estimate-fig}
to demonstrate the trend. One can of course examine the
role that the fidelity of the IMU plays in position and timing
estimation but as this plot demonstrate, there are more outliers
in the position estimation with only three SVs in view.
\begin{figure}[ht]
\begin{center}
\scalebox{0.8}{\includegraphics{eps/clockeps/3-sv-orbits-imu-estimate-enu.png}}
\caption{Rover position estimation in the ENU frame with 3 SVs in view and an IMU} \label{3-sv-rover-imu-estimate-fig}
\end{center}
\end{figure} 

This study can certainly be examined further to reduce the
number of SVs in view as shown in Figures~\ref{2-sv-rover-imu-estimate-fig} and
\ref{1-sv-rover-imu-estimate-fig} below.
\begin{figure}[ht]
\begin{center}
\scalebox{0.8}{\includegraphics{eps/clockeps/2-sv-orbits-imu-estimate-enu.png}}
\caption{Rover position estimation in the ENU frame with 2 SVs in view and an IMU} \label{2-sv-rover-imu-estimate-fig}
\end{center}
\end{figure} 
%-----
\begin{figure}[ht]
\begin{center}
\scalebox{0.8}{\includegraphics{eps/clockeps/1-sv-orbits-imu-estimate-enu.png}}
\caption{Rover position estimation in the ENU frame with one SV in view and an IMU} \label{1-sv-rover-imu-estimate-fig}
\end{center}
\end{figure} 

The situation with one SV in view of the rover would be
more manageable if we were not estimating the position over
the entire orbital period- the biases and noise
terms in the IMUs in general are integrated for velocity
and position estimation and without new measurements (say from SVs or
a base station) they will drift over time.


% We note in the meantime
% that the time estimation, even with one SV in view can 
% still be manageable as shown in Figure~\ref{}1-sv-rover-imu-time-estimate-fig

% \begin{figure}[t]
% \begin{center}
% \scalebox{0.8}{\includegraphics{eps/clockeps/rover-clock_1sv.png}}
% \caption{Rover clock bias estimation with one SV in view and an IMU} \label{1-sv-rover-imu-time-estimate-fig}
% \end{center}
% \end{figure} 

% In a nutshell you can build a model of the acceleration
% measurements available from the accelerometer.
% This measurement can then be combined with pseudorange
% measurement in the context of a Kalman filter.

% However, there are a number of ways to combine the
% Kuiper information with that of an accelerometer.


% We now consider a few scenarios to demonstrate some
% of the issues related to IMU integration. There are
% a few options to how to integrated 

% we adopt what is often referred to as tight integration
% of IMU to examine some of the 


% \subsection{Kuiper only velocity estimation}

% can we use the Kuiper directly for velocity estimation;
% this can be done with doppler and without and performance
% can be evaluated accordingly.

%\subsection{Example runs}

% the measured acceleration is actually 
% related to the linear acceleration,
% Corilolis acceleration, and gravitational acceleration- we will develop a more elaborate
% model when we we discuss aerial vehicles.

% once acceleration is obtained,
% we can then integrate to get velocity
% and then position- 

% \beas
% \bm a = \frac{d^2 \bm r}{dt^2} 
% \eeas

% and 

% \beas
% \bm v(t) = \int \bm a \, dt = a t + v_o
% \eeas 

% integration --- so have the clock bias corrected
% then use that for integration-
% loosely coupled and tighthy couopled


% The working of the accelerometer is
% as follows- and this is how it can be
% integrated with the Kuiper time stamping-
% we can also have an encoder on the wheels.

% The IMU works as follows--

% essentially it is a moving mass inside
% the chasis- so once we have three
% different d

% the output of the IMU 


% Another topic of interest is to combine time stamping scheme from Kuiper SVs
% with those from GPS for more precision positioning and time for Kuiper mobile units.

% % \subsection{Synchronization for Aerial Units}

% In this section, we consider synchronization problems for
% aerial units- the scenario of interest a two time tag
% setup as the vehicle flies over a CT-

% suppose we want to synchronize the timing between CT and aerial vehicle-
% the vehicle can use two time stamps 

% to correct with timing- we look at two time stamping scenario-

% this might be interesting to look at-- two way communication for synchronization




% \begin{itemize}
% \item using the GPS speed estimator
% \item using an accelerometer 
% \item digital compass
% \item rate gyros/IMU
% \item how to integrate GPS with IMUs
% \end{itemize}


% \section{Aerial Vehicle Positioning and Timing}

% We can can Kuiper SV time-stamped signals on aerial vehicles
% for positioning and timing- the baseline model can also
% be augmented with IMUs to better performance.

% \subsection{Without IMUs}


% \subsection{With IMUs}


% \section{Summary and Potential Next Topics}
% %
% In this chapter, we have examined how the Kuiper SV network
% can be used for positioning and time synchronization for
% mobile units. In the next chapter, we consider the
% same scenario for areial vehicles.

\subsection{Position/clock bias estimation for mobile units using time stamping and a base station}
%-----
Let us now go back GPS observables; these oberservables
are used to estimate position and clock bias of
the GPS receiver. These observables are constructed
differently but one of the step is to estimate the
signal emission time from GPS, namely, $t_e$;
we will discuss how these estimate are computed
subsequently in this chapter- and why for example, knowledge of the position
of the receiver leads to the sufficiency of
{\em one measurement} to obtain an estimate of the clock bias.

A mentioned above, mobile units can use propagation time
estimates from four SVs to compute their respective
range with respect to the satellites, and through these
measurements, obtain an estimate of its position
and clock offset. 
%----
And what is even more interesting, is
that we can use ``differential'' schemes,
similar to what is used in differential GPS (but
with respect to the SV signal),
or GPS-INS integration for mobile units,
but this time with time-stamping schemes.
%--------
Furthermore, analogous to the SV-CT sync, or the use of ISL,
we and utilize (extended) Kalman filtering
or network differencing to further improve the
synchronization performance. This will be done
both for the case that we have access to
position information, 
as well as when the position information has to be estimated.

How knowledge of the base's position,
and differencing schemes with a ``rover'' unit
be used for clock correction/positioning of the
mobile unit.


% \section{Velocity Estimation using Kuiper time signaling}


% \subsection{Positioning and Timing from LEO: Carrier tracking only}

% Along the lines of the reported work on Starlink;
% discuss in meeting. with a base station or other units.


% \subsection{Synchronization between Aerial and Ground Mobile Units}

% \subsection{Positioning and timing with GPS and Kuiper (Hybrid)}

% \subsection{Distributed Kalman filterin for position and clock estimation}

\subsection{Receiver Processing via Carrier Beat Frequency}

The basic ``observable'' in GPS receiver is a pseudorange,
a construct that has range, transmit times, transmit delays,
and processing delays all embedded; the pseudorange
is the range computed with all the errors-- hence the ``pseudo''!

There is actually two types of pseudoranges, one constructed
from code phase (difference) and one from the (carrier) beat frequency.
%--------
In a nutshell, once a particular GNSS signal has been acquired,
between $\mbox{SV}_i$ and $\mbox{CT}_k$,
the pseudrange (from code phase),
\beas
\mbox{pseudorange} = \mbox{true range} + c \; \mbox{(delays)}
\eeas
where $c$ is the speed of light. The ``delay'' term above is decomposed
and some of the terms can be accounted for in general, and some are
left as noise (more on this below). One term that is difficult to get ride of is
the receiver clock bias, particularly important since it multiplies the
speed of light in the pseudorange computation. 

A similar expression can be derived for the carrier phase
pseudorange (converting a phase to length by multiplying
by the wavelength) as,
\beas
\Phi_{i,k}= \rho_{ik} + c \; \mbox{(delays)} + \lambda \mbox{(initial phase offset)} + \lambda N_{i,k} - \epsilon_{i,k}
\eeas
where $N_{i,k}$ cycle ambiguity between the SV and CT;
the terms that we have ignored are the delay
uncertainties due to ionosphere and troposphere;
last term can capture for example delays due
to multipath.



The algorithmic sequence for GPS receiver software consists of,
\begin{enumerate}
	\item sample the channels that are tracking the GPS signal from a particular SV; we denote these times as $\tau_{\mbox{\tiny R}}(t_{r_i})$, where ``R'' is used generically for the receiver, and $t_{r_i}$ is the $i$th reception time.
	\item Based on the channel variables such as bit counts, number of C/A repetitions and code chips, $\tau_{{\mbox \tiny SV}_j}(t_{e_i})$ is determined; this is the signal emission time with resepct to the satellite clock. We will discuss this step in more detail.
	\item Using $\tau_{{\mbox \tiny SV}_j}(t_{e_i})$ and $\tau_{R}(t_{r_i})$, corrections sent by SV to the receiver, and other corrections included in the measurement, the transmission time $t_{e_i,r_i}= {t}_{r_i}- t_{e_i}$ is estimated; let us denote this by $\widehat{t}_{e_i,r_i}$.
	The pseudorange (for that channel) is now simply the speed of light multiplied by $\widehat{t}_{e_i,r_i}$.
\end{enumerate}
Once we have ``sufficient'' number of pseudoranges measurements,
we can then proceed to go after the position and the
clock offset for the receiver. 

Let us see how this actually works out
and in particular, how knowledge of the receiver position
would allow characterizing the receiver clock bias very efficiently.
The computation of the position becomes relevant when we 
subsequently consider mobile units.

Let the true range between the GPS satellite indexed as SV $i$ and the receiver
``R'' at time $t$ be,\footnote{The notation $\| \bm x \|$ denotes the 2-norm of the vector $\bm x$.}
\beas
\rho_{i,\mbox{\tiny R}}(t) = \| \bm r_{\mbox{\tiny SV}_i}(t) - \bm r_{\mbox{\tiny R}}(t)) \| ;
\eeas 
then the pseudorange can be written as,
\beas
\widehat{\rho}_{i, \mbox{\tiny R}}(t) = \rho_{i, \mbox{\tiny R}}(t) + c \; \delta t_{\mbox{\tiny R}} + e_{\mbox{\tiny R}}(t)
\eeas 
where $e_{\mbox{\tiny R}}(t)$ captures the stochastic errors (the sources of this error term consists of dispersive and non-dispersive atmospheric effects, multipath, etc.).

Note that in general the position vectors $\bm r_{\mbox{\tiny SV}_i}$ and 
$\bm r_{\mbox{\tiny R}}$ are often expressed in different coordinate frames;
let us for a moment assume that the relevant transformation has been
done to put them in the consistent coordinate frame and
\beas
\bm r_{\mbox{\tiny SV}_i}(t) = \leftm{c} \bm r^1_{\mbox{\tiny SV}_i}(t)\\[.05in]
										\bm r^2_{\mbox{\tiny SV}_i}(t)\\[.05in]
										\bm r^3_{\mbox{\tiny SV}_i}(t) \rightm
\quad \mbox{and} \quad
\bm r_{\mbox{\tiny R}_j}(t) = \leftm{c} \bm r^1_{\mbox{\tiny R}}(t)\\[.05in]
										 \bm r^2_{\mbox{\tiny R}}(t)\\[.05in]
										 \bm r^3_{\mbox{\tiny R}}(t) \rightm,
\eeas 
where the superscripts are the components of the corresponding
vector if that ``consistent'' coordinate frame.
%
This is important, as conversion between two coordinate frame one of which is time-varying, we need to use a consistent notion of time- which is also an unknown!

Let us for a moment assume that both position vectors have been
expressed in consistent coordinate frame; in this case, the
key algorithmic component of the GPS receiver is to use 
\beas
\widehat{\rho}_{\mbox{\tiny R}}^{j}(t), \quad j = 1, 2, \ldots
\eeas
to solve for $x_{\mbox{\tiny R}}(t)$, $y_{\mbox{\tiny R}}(t)$
$z_{\mbox{\tiny R}}(t)$ {\em and} $\delta t_{\mbox{\tiny R}}$.
In fact, using this pseudorange observable, the receiver formulates a least squares problems in four unknowns (three positions and one clock offset).

% \subsection{Position/clock bias estimation for mobile units using time stamping and beat frequency estimation}

% How $t_p$ can be used in conjunction beat frequency tracking
% to for clock/position estimation.

% \subsection{Position/clock bias estimation for mobile units using time stamping, beat frequency estimation, and inertial navigation}

% \begin{comment}
% Now let us consider the analogous construct with Kuiper. In this case,
% we do not necessary have a pseudorange construct; instead we 
% assume that $\tau_{\mbox{\tiny SV}}(t_e)$ is communicated to
% the receiver.

% the key difference if we have $t_e$ directly, ....

% Now suppose we 


% in the internal document it was stated that the SV $j$ sends
% INE-BCH (once per SDN cylce) provides SV ephemeris data,
% MAC subframe number, orbit information and GPS system time-
% we denote this last information by $\tau_{{\mbox \tiny SV}_j}(t_{e_i})$.
% For our purposed we assume that 
% \beas
% \tau_{{\mbox \tiny SV}_j}(t_{e_i}) = t_{e_i} \quad \quad \mbox{for all SV's $j$};
% \eeas
% but as we have discussed in this report, the discrepancy between 
% $\tau_{{\mbox \tiny SV}_j}$'s with each other and with $t$'s 
% can be measured using the network-bias and coherency that
% quantify {\em with and without} GPS 
% and even with {\em partial GPS} availability.

% What are the observables: in GPS case
% and Kuiper. What would a differential Kuiper 
% processing looks like-
% how can you use the communication link for navigation.

% why we can get away with one measurement for time?

% the difference between observables for GPS
% and observables for Kuiper.

% how differential GPS idea works and improves\\

% \end{comment}
% %
% \subsection{Carrier phase observable schemes}

% \subsection{Differencing schemes}
% \subsection{Differential methods for mobile units synchronization}

% \subsection{Cascade filtering}

% we can use the differential model to correct
% for the clock and position errors.
% We can use a SV network be used
% for positioning/timing of aerial vehicle that in turn can be used for position and timing of the ground vehicles.

% \subsection{Network-level synchronization}

% Can we use inter-mobile synchronization 
% (e.g., aerial to ground) be used to further improve synchronization performance.

% % ISL control system
% % spatial invariance for orbits and its consequences

% \begin{comment}
% \beas
% \bm q_{\mbox{\tiny R}}(t) = \leftm{c} \bm r_{\mbox{\tiny R}}(t) \\
% 						   \bm x_{\mbox{\tiny R}}	(t)
% 	   				      \rightm
% \eeas 

% will use superscripts to denote the coordinates- so for example,
% $\bm x^{1}_{\mbox{\tiny R}}(t)$ would be the clock bias for the receiver
% at time $t$.

% and based on the pseudorange measurements, we aim to estimate 

% \beas
% {\widehat{\bm q}}_{\mbox{\tiny R}}(t)}
% \eeas 
% where as before 
% \beas
% \bm x_i (t)= \leftm{c}
% \mbox{position for node $i$ at time $t$} \\
% \mbox{velocity for node $i$ at time $t$} \\
% \mbox{antenna orientation for for node $i$ at time $t$}\\
% \mbox{clock reading for for node $i$ at time $t$};
% \rightm
% \eeas 

% \beas
% A(t) \bm r_{\mbox{\tiny R}}(t) + \ones \, c \, \delta t_{\mbox{\tiny R}}(t) = \bm b_{\mbox{\tiny R}} (t)- \bm \epsilon_{\mbox{\tiny R}}(t)
% \eeas 


% \beas

% A(t) \bf r_{\mbox{\tiny R}} + \leftm{cccc} \ones & 0 & \cdots & 0 \\
% 										0     & \ones & \cdots & 0 \\
% 										\vdots & \vdots & \vdots & \vdots\\
% 									0     & 	0 \cdots & \ones \rightm

% 									= \leftm{c} c \delta t_{\mbox{\tiny R}} \rightm
% \eeas 

% how to use phase observations?

% Let $\ones = $

% \subsection{Differencing Schemes}

% differencing schemes are use to remove common errors

% rover and base setup: the base has knowledge of its position
% and shares carrier phase observables with the rover-
% the rover has no knowledge of its position but
% only the carrier phase observable.
% can shares 


% carrier phase 

% if the timing information is available ...

% we can also have a model of the receiver clock- if
% we don't then need to extract the nav. information somehow.



% \subsection{Comm for Nav}


% carrier phase observable with respect to the $i$th SV

% \beas
% \zeta(k) = \| xxxxx \| + a_i k D + b_i + c + \nu
% \eeas 

% in this case, we can solve for $a_i$'s and $b_i$ as well as 
% $r$ $D=10$ is the downsampling factor,


% This letter makes the following contributions. First, the Starlink signals are analyzed and a model suitable for carrier phase
% tracking is developed. Second, an adaptive KF-based tracking
% loop is developed where the measurement noise is updated
% based on a heuristic of the residuals. Third, a demonstration
% of the first carrier phase tracking and positioning results with
% real Starlink signals is presented, showing a horizontal position
% error of 7.7 m with six Starlink SVs.


% The adaptive KFbased carrier tracking algorithm is described below.


% ir interface in general, except for the channel frequencies
% and bandwidths. One cannot readily design a receiver to track
% Starlink signals with the aforementioned information only as
% a deeper understanding of the signals is needed. Softwaredefined radios (SDRs) come in handy in such situations, since
% they allow one to sample bands of the radio frequency spectrum. However, there are two main challenges for sampling
% Starlink signals: (i) the signals are transmitted in Ku/Kabands, which is beyond the carrier frequencies that most
% commercial SDRs can support, and (ii) the downlink channel
% bandwidths can be up to 240 MHz, which also surpasses the
% capabilities of current commercial SDRs. The first challenge
% can be resolved by using a mixer/downconverter between the
% antenna and the SDR. However, the sampling bandwidth can
% only be as high as the SDR allows. In general, opportunistic
% navigation frameworks do not require much information from
% the communication/navigation source (e.g., decoding telemetry
% or ephemeris data or synchronizing to a certain preamble).
% Therefore, the aim of the receiver is to exploit enough of the
% downlink signal to be able produce raw navigation observables
% (e.g., Doppler and carrier phase). Fortunately, a look at the
% FFT of the downlink signal at 11.325 GHz carrier frequency
% 0018-9251 (c) 2021 IEEE. Personal use is permitted, but republication/redistribution requires IEEE permission. See http://www.ieee.org/publications_standards/publications/rights/index.html for more information.
% This article has been accepted for publication in a future issue of this journal, but has not been fully edited. Content may change prior to final publication. Citation information: DOI 10.1109/TAES.2021.3113880, IEEE
% Transactions on Aerospace and Electronic Systems
% and sampling bandwidth of 2.5 MHz shows nine “carrier
% peaks,” as shown in Fig. 1(a). Furthermore, the waterfall plot
% in Fig. 1(b) shows that these carrier peaks vary as the Doppler
% frequency over an 80-second interval. The Doppler frequency
% was predicted using two-line element (TLE) files

% \end{comment}

% \section{Synchronization with the mobile units}

% \subsection{GPS to Mobile Units}



% \begin{itemize}
% \item do we have to worry about multipath for mobile units? can we use mobile units in concert to eliminate the effect of multipath?
% \end{itemize}

\section{Kuiper Aerial Vehicle Network}
In this section, we discuss the use of Kuiper Aerial Vehicles (KAVs) network for positioning and timing. In particular, we consider how to use
KAVs as ``anchors" for Kuiper ground mobile units or IoT devices, 
or augmenting timing and navigation signals from 
GPS or Kuiper satellites to improve timing and position
estimation for other Kuiper nodes.

Kuiper aerial vehicles provide a unique opportunity to
augment the baseline Kuiper network. This augmentation
has multiple facets. For example, long endurance medium altitude
aerial vehicles (operating altitudes between 3-7 km for 24-48 hours)
or long endurance high altitude aerial vehicles (operating altitude
of 20 km for months) provide an unprecedented opportunity for
creating ad hoc communication or navigation solutions to customers
on the ground and sea. Such vehicles can also be used for augmenting
existing networks or such as Kuiper SV/GW/CT (for communication and navigation)
as well as navigation networks such as GPS/GNSS.

At lower altitudes, in addition to augmenting
existing satellite or ground networks, KAVs can provide
other types of functionalities. For example, the KAV clock can be synchronized
with ``Kuiper time" as it flies over a customer terminal or a
base station, and subsequently be used ``relay" that information
to ground mobile units or ships. KAVs can
also be used to provide network connectivity during 
disruptions in the baseline satellite networks.

Let us now examine a few scenarios of interest.

\begin{itemize}

\item Consider the situation where there has been a 
GPS disruption at a ground mobile unit or a region of interest.
As discussed in the previous chapter, in such scenarios,
one can use the time stamped signals from Kuiper SVs for positioning and time
synchronization of the ground units. This setup would be more efficient
and accurate when we have at least four Kuiper SVs in view 
(without using any other ``side information'' like knowledge of the
terrain). As we also showed in the previous chapter, one can use the
Kuiper SV time-stamped signals in conjunction with inertial measurement units
to improve the positioning and timing performance when the ground
rover does not have access to four SVs during the clock steering 
and position/velocity updates. However, the above scheme can
only work if the rover's antenna can receive the required time-stamped signal
from the SVs. For this to happen, the antenna on the ground vehicle
should be pointing toward the transmitter's antenna, implicitly
suggesting that the rover should know where it is with respect
to the satellites on the earth's surface. 

\begin{figure}[h]
\begin{center}
\includegraphics[width=.5\textwidth]{eps/clockeps/mobility-kav-base.jpg}
\caption{Synchronization for aerial and Kuiper base using two time stamped signaling}\label{kav-base}
\end{center}
\end{figure}

If the rover (or a Kuiper boat) does have any information about its location nor
where in the sky to search for Kuiper satellites, then a
solution strategy, based on a KAV node can be adopted
as follows. The KAV first identifies a Kuiper ground vehicle 
(say from a tag or identifiable feature) using its onboard camera.
If the KAV camera is gimbaled, it can scan areas on the ground
where the likelihood of the ground vehicle finding a satellite is too
low. Once the camera of a KAV sees (pixels on the camera image plane!)
the ground vehicle, it can
estimate the ground vehicle's position from the projection of
the vehicle in the camera frame. There is some geometry 
and coordinate transformation involved in this process, but this can be done.
We can even use multiple KAVs to improve the rover's position/antenna
orientation estimation by running a distributed estimation algorithm
across the KAV network flying over the ground receiver.

\item Note that in the scenario above, the KAV needs to know
its altitude with a use of an altimeter. Alternatively, we can
assume that KAV can estimate its position and orientation, and
as such, can estimate its states using GPS or Kuiper SV network.

\item Now suppose that KAV identifies the Kuiper mobile unit
that does not have a GPS access (may be there is status indicator on
the mobile unit). Furthermore, assume 
that the mobile unit's antenna is not pointing in the right
direction. If the KAV's altitude is low, then it can 
use its camera and some ML, to estimate the orientation of
the antenna and then choose a trajectory to fly in the
footprint of rover's antenna. We can then develop a steering 
algorithm on the KAV's antenna to keep the rover in its
foot print for the time duration needed to perform
clock synchronization and facilitate a navigation solution for 
the rover. This information can subsequently be used by the rover,
in order to scan the correct portion of the sky for SV signals
and subsequent synchronization and positioning. The KAV initialization
can also be used by the rover's onboard IMU for state propagation.

\item In the scenario above, one is using a low altitude
KAV with a gimbaled camera to resolve for the relative position
(vector) of the ground rover with respect to the
aerial vehicle. In a nutshell, we are using the onboard camera
onboard- in addition to some geometric relations-
to estimate the 3D relative position vector and subsequently achieve
time synchronization. Now suppose that instead, we use
signal processing to estimate the delay
in the signal propagation from KAV to the rover. 
Then having a model for this delay,
and the positions of at least four KAVs, would allow the
rover to estimate its position in the world coordinate.
If KAVs' velocity information is also available to the rover,
then doppler can be used to estimate the velocity of
the rover. Hence, in essence, one can create a 
{\em KAV anchor network} that can be used for positioning and timing of 
other vehicles in the network as needed. Since the accuracy of the
navigation and timing is a function of the geometry of the
anchor units, we can then choose the KAV's configuration (adaptively!)
for a given region that minimizes the average estimation error for
both timing and navigation performance.

\item Lastly, KAV can act both as relays as well as augmentation
units for the comm network when needed- creating a adaptive "node"
that can effectively be used during disruptions, e.g., 
disaster recovery.

\item And of course, KAV can also be used to augment Kuiper SV or GPS signals
for the rover positioning and timing, including facilitating differential
processing for higher accuracy navigation and timing.
\end{itemize}

\begin{figure}
\begin{center}
\includegraphics[width=.5\textwidth]{eps/clockeps/Kuiper-mobile-KAV.jpg}
\caption{Synchronization for aerial and CT using two time stamped signaling}\label{mobility-kav}
\end{center}
\end{figure}

\begin{figure}
\begin{center}
\includegraphics[width=.5\textwidth]{eps/clockeps/Kuiper-mobile-KAV2.jpg}
\caption{Synchronization for aerial and CT using two time stamped signaling}\label{mobility-kav2}
\end{center}
\end{figure}


\section{Localizing Ground Mobile Units from the Air using an Onboard Camera}
In this section, we examine how to use an
onboard camera on a KAV for positioning and time synchronization 
of a mobile ground rover. This setup provides the means of
{\em initializing} the subsequent positioning and timing,
say through delay estimation, as well as air-ground communication. 
Hence, the setup can be referred
to as localization without relying on communication, or
localization in support of communication!
%
Let us explain the setup. It is assumed that KAV has synchronized its
time with Kuiper base and ``in-tune'' with 
the system time; it also has self navigation information.
Furthermore, we assume that
KAV has identified a rover that requires
time/navigation information- this can be 
either through an SOS message as Pengfei suggested
or having KAV(s) cover an area and looking for a status
light on the rover that has lost its navigation/timing information.

Suppose that the origins of the camera frame and that of the gimbal,
onboard the KAV, coincide with the center of mass of the KAV. Suppose that
we also have the body frame for the KAV.
Now we monitor the rotation that is needed to go from the
body frame to the gimbal and the camera frames. These coordinate
transformations involve a few angles, including the so-called
gimbal-azimuth $\alpha_{z}$ and the gimbal elevation $\alpha_{z}$ 
angles. % These are shown in the figure below. 
Using these angles, we can find the rotation matrix needed to go 
from the body to the gimbal frame; hence using these angles and
parameterize the direction cosine matrix ${\cal C}_{b}^{g}$.
%
We also have a transformation to go from the gimbal frame to the camera
frame that we have to keep track of; the computer vision folks
generally set these coordinate frames in such a way that
the corresponding rotation matrix is of the form,
\beas
{\cal C}_{g}^{c} = \leftm{ccc} 0 & 1 & 1 \\ 0 & 0 & 1\\ 1 & 0 & 0 \rightm.
\eeas
Hence, so far, we can find the necessary coordinate
transformations between the gimbal, camera, and the body frames of KAV
in terms of these two angles.

Now, let us go back to the camera frame. A typical camera
geometry is as follows. The key parameters are the focal 
length in units of pixels $f$ and its conversion to
units of meters $pf$ ($p$ is a scaling factor).

\begin{figure}[h]
\begin{center}
\includegraphics[width=.5\textwidth]{eps/clockeps/camera-frame.jpg}
\caption{The camera geometry}\label{camera-frame}
\end{center}
\end{figure}
Let us assume that the pixel array is square for now consisting also of square pixels; 
let the (pixel) width of the pixel array be $M$. The camera focal length 
and the array width are the related through the
half angle of the camera field of view;
in particular, 
\beas
\tan \frac{\nu}{2} = \frac{M}{2 f};
\eeas 
so in essence we know $f$ and $pf$ if we know the camera field of view and the array (pixel) width.

Now comes the interesting (yet simple) relationship.
In relation to Figure~\ref{camera-frame}, 
suppose the pixel coordinates 
in the image plane of the
KAV camera corresponding to rover is 
$(\epsilon^x_{\mbox{\tiny R}},\epsilon^y_{\mbox{\tiny R}})$; in the 3D version, this coordinate is
really
\beas
(\epsilon^x_{\mbox{\tiny R}},\epsilon^y_{\mbox{\tiny R}}, pf).
\eeas 
Therefore the distance from the camera frame
to the pixel location is 
\beas
F= \sqrt{f^2 + (\epsilon^x_{\mbox{\tiny R}})^2 + (\epsilon^y_{\mbox{\tiny R}})^2  }.
\eeas 
We use some geometry (similar triangles) to realize that,
\beas
\bm r^c_{\mbox{\tiny R/KAV}} = \frac{d}{F} 
\leftm{c} \epsilon^x_{\mbox{\tiny R}}\\
\epsilon^y_{\mbox{\tiny R}}\\
f \rightm
\eeas 
where $d$ is the distance between the rover and
its projection on the image plane.

This means that although the relative
vector from the KAV to the rover can not be
uniquely determined, its ``direction'' can!
That is,
\beas
\widehat{\bm r}^c_{\mbox{\tiny R/KAV}} = \frac{1}{\sqrt{f^2 + (\epsilon^x_{\mbox{\tiny R}})^2 + (\epsilon^y_{\mbox{\tiny R}})^2  }} 
\leftm{c} \epsilon^x_{\mbox{\tiny R}}\\
\epsilon^y_{\mbox{\tiny R}}\\
f \rightm;
\eeas 
we are adopting the convention of referring to the normalized version of any vector $\bm v$ 
as $\widehat{\bm v}$.

Since we can find the direction to the rover, all we need know now is the range to the rover. 
That is, once we know this range, the we can obtain
${\bm r}^c_{\mbox{\tiny R/KAV}}$ and convert
that to ${\bm r}^i_{\mbox{\tiny R/KAV}}$
that can be combined with 
${\bm r}^i_{\mbox{\tiny KAV}}$ to obtain
${\bm r}^i_{\mbox{\tiny R}}$, i.e., 
\beas
\bm r_{\mbox{\tiny R}/i}  = \bm r_{\mbox{\tiny KAV}/i} + {\cal C}_{b}^i \, {\cal C}_{g}^b \, {\cal C}_{c}^g \, {\bm r}^c_{\mbox{\tiny R/KAV}};
\eeas 
recall that if the superscript is not included, the vector is being resolved in the inertial frame.
%
Now suppose that we know the altitude of the KAV denoted by $h$. 
Then the geometry of Figure~\ref{mobility-kav}
with a flat earth approximation, implies
that, 
\beas
d = \frac{h}{\cos \phi}
\eeas 
where $\phi$ is the angle between 
${\bm r}^c_{\mbox{\tiny R/KAV}}$,
converted to the inertial frame,
and $k$ axis of the inertial frame. 
This completes the relation location 
determination of the rover 
by the gimbaled camera on the KAV.
%
The resolution of such a localization
depends camera resolution and as knowledge
of the various angles and the height
of the KAV. The basics are there, but of course,
there are a number of ways to improve this
estimate, including using multiple KAVs
to localize a rover and then exchange information
to reduce the localization uncertainty; another
possibility is to use a filter to update the rover's location 
using a time series of camera images. The localization can
also be improved by strategically maneuver a low altitude KAV 
so that the uncertainty in the measurement geometry is reduced.

Once the relative position of the rover is determined, the 
KAV can perform a ``initialization'' maneuver followed by an 
initialization protocol for the rover, to provide the information
to the rover to updates its position, thus facilitating subsequent 
communication with the KAV or KSV networks.
% %

% \subsection{Keeping the rover in the KAV camera field of view}

% \subsection{Simultaneous localization and time synchronization}

% \subsection{Air-ground communication}

% \section{KAV network for positioning and timing}


% \section{KAV-SV network for positioning and timing}

% Once we know the relative position of the rover - and either estimate the pose of the antenna
% from the KAV camera or know the default orientation (for example, rover antenna should
% go) ... check what type of camera and then see the rover antenna geometry that is needed
% to intercept the beam.

% \section*{Notes and References}




% \begin{figure}[ht]
% \begin{center}
% \includegraphics[width=.8\textwidth]{eps/clockeps/mobility1.jpg}
% \caption{Synchronization for aerial and CT using two time stamped signaling}\label{mobility1}
% \end{center}
% \end{figure}


\begin{figure}[h]
\begin{center}
\includegraphics[width=.4\textwidth]{eps/clockeps/mobility2.jpg}
\caption{Positioning using mixed time stamped signaling from aerial units and Kuiper SV}\label{mobility}
\end{center}
\end{figure}



\section*{Notes and References}
In this chapter, we examine how to using Kuiper time-stamping
not only for time-synchronization but also for positioning,
adopting a one-way time transfer protocol. In this case,
timing and positioning becomes coupled, as such the
time-stamped signals are used in conjunction for resolve
for position and clock offset of a ground mobile vehicle.
When enough SV time-stamped signals are not available, the
vehicle was augmented with an inertial measurement unit.
It is show that the Kuiper SV time-stamped signaling 
in conjunction with the vehicle inertial measurement provide
a rather robust mechanism for time synchronization and
position estimation. The positioning and timing protocols
discussed in this chapter are part of the broader 
discipline of positioning, navigation, and timing (PNT).
This is an active area of research currently,
as there are a number of extensions that can be 
examined in the context of Kuiper network. One prime
example in this direction is the use of Kuiper aerial vehicles, both in the backup and augmentation modes for
the baseline Kuiper network.

% \begin{figure}[h]
% \begin{center}
% \includegraphics[width=.4\textwidth]{eps/clockeps/mobility3.jpg}
% \caption{Positioning using mixed time stamped signaling from aerial units and Kuiper SV}\label{mobility}
% \end{center}
% \end{figure}


% \section*{Notes and References}
% \chapter{Control system for ISL Instrument}


% \chapter{IoT Synchronization/Indoor Positioning \& Timing}

% latency issues for critical systems (industrial, autonomy), energy constraints/synchronization performance. Also interacting with the environment
% and how learning/data driven methods can
% be used to optimize energy and lower the
% communication overhead to achieve 
% the desired levels of synchronization.


% \section{Indoor Positioning and Timing}

% \section{Energy-aware synchronization}

% \section*{Notes and References}
% %learning/data driven?

% \part{Part III: System-level Kuiper Network Synchronization}
%\include{constellations-orbital}

%\include{iterative-design}
% %-----
% \part{Part IV: Kuiper Aerial Vehicles}

% \part{Part V: Internet of Things}

%\include{state-coordination}

%\include{network-resilience}
%\part{Summary and Future Works}
\include{summary-future-works}
%\part{Appendix}
%-------
%\include{conference_paper}
% \include{matrices-probability-optimization}
% \include{gps-theory}
%\include{interference}
%\include{other-topics}
%\include{comments_from_kuiper_team}
%------

% \include{extra}
%\part{Code}
\bibliographystyle{plain}
\bibliography{references}
\end{document}

% 
% other topics could include DTN performance
% lunar communication and navigation
% as well as the use of spatial invariance and symmetry
% properties of the Kuiper network
